\documentclass[12pt,a4paper]{article}

% --- Преамбула ---
\usepackage[utf8]{inputenc}
\usepackage[T1]{fontenc}
\usepackage[russian]{babel}
\usepackage{amsmath,amssymb}
\usepackage{hyperref}
\usepackage[margin=2.5cm]{geometry}
\usepackage{enumitem}
\usepackage{setspace}
\onehalfspacing

\title{\textbf{Теория Мод}\\
\large Полное собрание глав (1--13)}
\author{}
\date{}

\begin{document}
\maketitle
\tableofcontents
\newpage

% ==================== Глава 1 ====================
\section{Теория Мод: Эмерджентное Пространство-Время, Квантовая Механика и Единая Картина Фундаментальных Взаимодействий}

\begin{abstract}
Предложена теория, в которой пространство-время, квантовая механика, гравитация и поля Стандартной Модели возникают как эмерджентные явления из единой фундаментальной структуры --- сети мод. Мода определяется как двумерная поверхность между двумя частицами, несущая продольную (счёт посредников) и поперечную (спектральную) структуру. Расстояние, масса, энергия, спин и заряд выводятся из геометрии спектров, а не постулируются. Причинность первична и формулируется как двустороннее условие существования: событие-причина~$A$ существует тогда и только тогда, когда существует событие-следствие~$B$. Время не имеет смысла для Мультиверса в целом, но определено как причинная последовательность внутри отдельных Вселенных. Из аксиоматики выводятся уравнения Эйнштейна, правило Борна, калибровочные симметрии, объяснения тёмной материи, тёмной энергии, барионной асимметрии и квантования. Предложены экспериментальные предсказания и явный алгоритм вычисления наблюдаемых величин через спектральные корреляции. Построен мост к стандартной квантовой механике: эффективная пси-функция взаимодействия двух частиц выводится как непрерывный предел фундаментальной суммы по дискретному пути.

\bigskip
\noindent\textbf{Ключевые слова:} эмерджентное пространство-время, моды, спектры, причинность, квантовая гравитация, тёмная материя, тёмная энергия.
\end{abstract}

\subsection{Введение}
Современная физика сталкивается с рядом фундаментальных проблем: противоречие между квантовой механикой и общей теорией относительности (ОТО), природа тёмной энергии и ускоренного расширения Вселенной, механизм квантовой запутанности и её масштабирование до макроскопических систем, происхождение пространства-времени (Лоренц-инвариантность, размерность $3+1$). Цель данной работы --- построить единую модель, разрешающую эти проблемы на основе концепции многомерных мод. В отличие от стандартных подходов, пространство-время здесь не фундаментально, а возникает как статистический предел дискретной сети спектральных корреляций.

\subsection{Часть I. Определения (Аксиоматика)}

\subsubsection{Определение 1. Частица}
Частица $P_i$ --- фундаментальный узел сети. Каждая частица обладает \textbf{собственным спектром} $S_i$ --- функцией, описывающей её внутреннее состояние. Все частицы являются составными: внутренняя структура частицы --- это ансамбль суб-частиц, чьи спектры интерферируют, формируя результирующий спектр $S_i$. Спектр $S_i$ не зависит от какого-либо внешнего времени; частица существует как целостный объект в сети.

\subsubsection{Определение 2. Мода}
Для каждой упорядоченной пары различных частиц $(P_i, P_j)$ существует \textbf{мода} $M_{i \to j}$. Это отображение из дискретного причинного пути в скалярное поле.

Пусть $\Gamma_{ij}$ --- дискретный путь (цепочка) от $P_i$ до $P_j$, существование которого обусловлено причинным порядком и спектральной связностью. Путь записывается как упорядоченное множество:
\begin{equation}
    \Gamma_{ij} = (P_i = P_{k_0}, P_{k_1}, P_{k_2}, \dots, P_{k_d} = P_j)
\end{equation}
где $d = |\Gamma_{ij}| - 1$ --- длина пути в числе шагов (число промежуточных узлов).

Тогда мода --- это функция на этом пути:
\begin{equation}
    M_{i \to j}: \Gamma_{ij} \to \mathbb{R}
\end{equation}
сопоставляющая каждой точке пути $P_{k_m}$ вещественное число $F_{i \to j}(m)$ --- значение поперечного спектра в этой точке.

\textbf{Структура моды:}
\begin{itemize}
    \item \textbf{Носитель (продольная структура):} сам дискретный путь $\Gamma_{ij}$ с индексом шага $m = 0, 1, \dots, d$.
    \item \textbf{Спектр (поперечная структура):} вектор значений $(F(0), F(1), \dots, F(d))$.
\end{itemize}

Каждая частица $P_i$ является источником набора мод $\mathcal{M}_i = \{M_{i \to j} \mid \forall j \neq i\}$ ко всем остальным частицам Вселенной. Состояние частицы --- суперпозиция всех её мод:
\begin{equation}
    S_i = \bigoplus_{j \neq i} M_{i \to j}
\end{equation}

Значение спектра на конкретной моде есть проекция собственного спектра частицы в данном направлении:
\begin{equation}
    F_{i \to j}(m) = \text{Proj}_{\angle(i \to j)} S_i
\end{equation}

\textbf{Принцип компенсации:} сумма значений спектра по всему пути стремится к нулю:
\begin{equation}
    \sum_{m=0}^{d} F_{i \to j}(m) \approx 0
\end{equation}
Мода в целом нейтральна, но локально асимметрична, что и порождает взаимодействие.

\subsubsection{Определение 3. Спектральная связность и критерий существования пути}
Не для всякой пары частиц путь $\Gamma_{ij}$ и мода $M_{i \to j}$ существуют. Условие существования --- \textbf{спектральная связность}:

Путь $\Gamma_{ij}$ существует $\iff$ корреляция между проекциями собственных спектров в направлении друг друга отлична от нуля:
\begin{equation}
    \langle F_{i \to j} \cdot F_{j \to i}^{\text{rev}} \rangle \neq 0
\end{equation}
где $F_{j \to i}^{\text{rev}}$ --- развёрнутый спектр встречной моды: $F_{j \to i}^{\text{rev}}(m) = F_{j \to i}(d-m)$.

Если корреляция равна нулю, частицы не взаимодействуют, и расстояние между ними не определено (или считается бесконечным).

\subsubsection{Определение 4. Расстояние}
\textbf{Расстояние} $d(P_i, P_j)$ между двумя частицами, для которых существует спектральная связность, определяется через понятие \textbf{причинно-спектрального пути}.

\textbf{Причинно-спектральный путь} $\Gamma_{ij}$ от $P_i$ к $P_j$ --- это упорядоченная цепочка частиц, удовлетворяющая двум условиям:
\begin{enumerate}[label=(\arabic*)]
    \item \textbf{Причинная последовательность:} $P_{k_m} \prec P_{k_{m+1}}$ для всех $m$ (см.~Определение~6).
    \item \textbf{Спектральная компенсация:} каждый шаг $P_{k_m} \to P_{k_{m+1}}$ уменьшает локальную нескомпенсированность.
\end{enumerate}

\textbf{Расстояние} --- это число промежуточных узлов в \textit{кратчайшем} (по числу шагов) причинно-спектральном пути:
\begin{equation}
    \boxed{d(P_i, P_j) = \min_{\Gamma} |\{P_k \in \Gamma \mid P_k \neq P_i,\; P_k \neq P_j\}|}
\end{equation}

Если существует несколько путей одинаковой минимальной длины, выбирается путь с наименьшей суммарной нескомпенсированностью.

\textbf{Критерий «между»:} Частица $P_k$ находится \textit{между} $P_i$ и $P_j$ тогда и только тогда, когда она принадлежит кратчайшему причинно-спектральному пути $\Gamma_{ij}$.

Это определение не рекурсивно: путь $\Gamma_{ij}$ строится из условия причинного порядка и минимизации спектральной нескомпенсированности на каждом шаге, без отсылки к заранее известному расстоянию.

\subsubsection{Определение 5. Взаимодействие}
Сила и тип взаимодействия между $P_i$ и $P_j$ определяются корреляцией двух встречных мод на кратчайшем пути $\Gamma_{ij}$ длины $d$:
\begin{equation}
    V_{ij} = f\left( \sum_{m=0}^{d} F_{i \to j}(m) \cdot F_{j \to i}(d-m) \right)
\end{equation}

Типы взаимодействий:
\begin{itemize}
    \item \textbf{Притяжение:} сумма близка к нулю (пики одного приходятся на провалы другого).
    \item \textbf{Отталкивание:} $|V_{ij}| \gg 0$ (пики совпадают, сумма велика).
    \item \textbf{Нейтральность:} $V_{ij} \approx 0$ (спектры некоррелированы).
\end{itemize}

\subsubsection{Определение 6. Причинность}
Причинность первична и является отношением частичного порядка $\prec$ на множестве всех частиц, удовлетворяющим глобальному условию самосогласованности:

\begin{quote}
\textbf{Событие-причина $A$ существует тогда и только тогда, когда существует событие-следствие $B$.}
\end{quote}
\begin{equation}
    \exists A \iff \exists B
\end{equation}

Это условие накладывается на всю сеть мод и определяет, какие причинные пути $\Gamma$ являются допустимыми. Фотон, который никогда ни с чем не провзаимодействует в будущем, не может быть испущен в прошлом.

\subsubsection{Определение 7. Время}
\textbf{Глобальное время (Мультиверс):} Для сети мод как целого время не определено. Вся совокупность частиц и мод существует как статичная структура. Прошлое и будущее --- не разные моменты, а разные области сети.

\textbf{Внутреннее время (Вселенная):} Внутри связной области сети (Вселенной) причинный порядок $\prec$ индуцирует параметр $t$, который воспринимается внутренними наблюдателями как время. Этот параметр:
\begin{itemize}
    \item не является фундаментальной координатой частицы;
    \item возникает как эмерджентное свойство причинной топологии;
    \item не обязан быть глобально линейным (допускает петли и ветвления, если они удовлетворяют $\exists A \iff \exists B$).
\end{itemize}

\textbf{Стрела времени:} Точка настоящего всегда неравновесна, поскольку прошлое и будущее бесконечны и не полностью симметричны по структуре. Это создаёт направленность причинной последовательности.

\subsubsection{Определение 8. Энергия}
Энергия не является самостоятельной субстанцией. Это количественная мера нескомпенсированности поперечных спектров в данной области сети:
\begin{equation}
    E(\mathcal{R}) = \sum_{(i,j) \in \mathcal{R}} \left| \langle F_{i \to j} \cdot F_{j \to i}^{\text{rev}} \rangle \right|
\end{equation}

Энергия не сохраняется глобально, но в пределе изолированной системы, где суммарная нескомпенсированность постоянна, восстанавливается стандартный закон сохранения. Рождение частицы-посредника --- способ превратить спектральный конфликт в расстояние, уменьшив локальную энергию.

\subsubsection{Определение 9. Масса}
Масса --- степень размытия пиков и провалов поперечного спектра частицы. Два источника размытия:
\begin{enumerate}[label=\arabic*.]
    \item \textbf{Внутренний состав:} Наложение мод внутри частицы (все частицы составные). Чем сложнее внутренняя структура, тем сильнее размыт результирующий спектр.
    \item \textbf{Внешнее сближение:} Интерференция с модами соседних частиц при сильном сближении.
\end{enumerate}

Формально:
\begin{equation}
    m_i \propto \Delta S_i + \sum_{j \in \text{neighbours}(i)} \kappa(d_{ij}) \cdot \Delta S_j
\end{equation}
где $\Delta S_i$ --- внутреннее размытие, $\kappa$ --- функция затухания с расстоянием.

Масса не постоянна и зависит от окружения. Безмассовые частицы (фотон) имеют идеально острый спектр без размытия. Максимальная масса (чёрная дыра) --- спектр, полностью размытый до однородного фона.

\subsubsection{Определение 10. Спин}
Спин --- внутренний градиент поперечного спектра вдоль моды. Это паттерн чередования ориентации пиков/провалов по продольной координате моды.

Формально:
\begin{equation}
    \sigma = (\sigma_0, \sigma_1, \dots), \quad \sigma_x = \text{sign}(F(x)) \in \{+, -\}
\end{equation}

\begin{itemize}
    \item \textbf{Спин 1/2 (фермионы):} $+ - + - \dots$ или $- + - + \dots$
    \item \textbf{Спин 0 (бозоны):} Синфазное чередование.
    \item \textbf{Спин 1:} Период чередования 1.
    \item \textbf{Спин 2 (гравитон):} Более сложный тензорный паттерн.
\end{itemize}

\textbf{Запрет Паули:} Два фермиона не могут занимать одно состояние, так как их спиновые последовательности должны быть различны --- иначе их моды были бы неразличимы.

\subsubsection{Определение 11. Античастица}
Античастица $\bar{P}$ --- конфигурация с инвертированным поперечным спектром:
\begin{equation}
    S_{\bar{P}} = -S_P
\end{equation}
Все пики становятся провалами, и наоборот. Спиновая последовательность также инвертируется. При наложении мод частицы и античастицы их спектры полностью компенсируются (аннигиляция).

\subsection{Часть II. Теоремы}

\subsubsection{Теорема 1. Эмерджентность Пространства-Времени}
Из дискретной сети мод с расстоянием, определённым как число посредников, в пределе больших $N$ возникает непрерывное псевдориманово многообразие с лоренцевой метрикой.

Процедура усреднения (coarse-graining) сопоставляет частицам координаты $x^\mu$ так, что $d(P_i, P_j)$ аппроксимируется геодезическим расстоянием $\int \sqrt{g_{\mu\nu} \dot{x}^\mu \dot{x}^\nu} d\lambda$. Причинный порядок $\prec$ индуцирует световой конус и лоренцеву сигнатуру.

Уравнения Эйнштейна возникают как термодинамическое уравнение состояния сети при максимизации энтропии связности:
\begin{equation}
    R_{\mu\nu} - \frac{1}{2} g_{\mu\nu} R = 8\pi G \, T_{\mu\nu}
\end{equation}
где $T_{\mu\nu}$ --- макроскопическое описание нескомпенсированности спектров.

\subsubsection{Теорема 2. Квантовая Вероятность и Правило Борна}
Амплитуда перехода системы между состояниями $\alpha$ и $\beta$ определяется интегралом перекрытия всех участвующих мод:
\begin{equation}
    \mathcal{A}(\alpha \to \beta) = \prod_{(i,j)} \langle M_{i \to j}^\alpha \mid M_{j \to i}^\beta \rangle_{\text{спектр}}
\end{equation}

Правило Борна $P \propto |\mathcal{A}|^2$ --- следствие требования инвариантности вероятности относительно выбора глобальной фазы спектров.

Измерение --- акт резонансной синхронизации спектров прибора и системы. Коллапс --- схлопывание в одну из возможных корреляционных пар, для которой выполняется глобальное условие $\exists A \iff \exists B$.

\subsubsection{Теорема 3. Тёмная Энергия}
Ускоренное расширение Вселенной --- следствие глобальной нескомпенсированности спектров, требующей непрерывного рождения новых частиц-посредников.

Плотность тёмной энергии:
\begin{equation}
    \rho_\Lambda \propto \frac{\dot{N}}{N}
\end{equation}
где $N(t)$ --- число частиц в причинно-связанном объёме. Уравнение состояния: $p_\Lambda = -\rho_\Lambda$ (отрицательное давление, так как рождение посредника увеличивает расстояние, а не «расталкивает» частицы).

В уравнениях ОТО это даёт космологическую постоянную:
\begin{equation}
    \Lambda \propto \langle \dot{N}/N \rangle
\end{equation}

\subsubsection{Теорема 4. Замедленная Декогеренция Макроструктур}
Две макроскопические системы с идентичной топологией спектров сохраняют взаимную квантовую корреляцию на временах, экспоненциально больших, чем время декогеренции отдельных частиц.

Декогеренция --- расхождение спектров под воздействием шума от мод окружения. Топологическая защита (масса и спин) создаёт энергетический барьер $E_{\text{барьер}}$ против изменения топологического инварианта $T$. Вероятность декогеренции:
\begin{equation}
    P_{\text{decohere}} \sim \exp\left(-\frac{E_{\text{барьер}}}{E_{\text{шум}}}\right)
\end{equation}

Для идентичных копий шум действует скоррелированно, и время декогеренции:
\begin{equation}
    \tau_{\text{decohere}} \sim \tau_0 \cdot \exp(\alpha N)
\end{equation}
где $N$ --- число частиц в системе, $\alpha$ --- коэффициент топологической жёсткости.

\textbf{Экспериментальное предсказание:} Два топологически идентичных фрактальных кристалла должны проявлять скоррелированные флуктуации с задержкой $\Delta t \ll L/c$, где $L$ --- расстояние между ними.

\subsubsection{Теорема 5. Калибровочные Симметрии}
SU(3), SU(2), U(1) --- классификация допустимых преобразований внутренней структуры спектра, не меняющих корреляции $\langle F_{i \to j} \cdot F_{j \to i}^{\text{rev}} \rangle$.

\begin{itemize}
    \item \textbf{U(1):} Вращение фазы в одномерном внутреннем пространстве (электромагнетизм).
    \item \textbf{SU(2):} Унитарные преобразования с $\det = 1$ в двумерном пространстве (слабое взаимодействие).
    \item \textbf{SU(3):} Унитарные преобразования с $\det = 1$ в трёхмерном пространстве (сильное взаимодействие).
\end{itemize}

Размерность внутреннего пространства спектра определяет калибровочную группу.

\subsubsection{Теорема 6. Квантовая Запутанность}
Запутанность --- проявление того факта, что две частицы всегда уже связаны парной модой $\mathcal{M}_{AB}$, и их спектры образуют единый, неразделимый корреляционный объём.

Совместный спектр $\mathcal{S}_{AB}$ не факторизуется. Расстояние (число посредников) не влияет на поперечную корреляцию, поэтому запутанность не ослабевает с расстоянием.

Измерение частицы $A$ изменяет её спектр $S_A$, что требует мгновенного изменения $S_B$ для сохранения согласованности $\mathcal{M}_{AB}$. Это коллапс на расстоянии, не требующий передачи сигнала.

\subsubsection{Теорема 7. Спектр Реликтового Излучения}
Реликтовое излучение --- прямое наблюдение нескомпенсированности спектров ранней Вселенной на поверхности последнего рассеяния.

Анизотропии температуры:
\begin{equation}
    \frac{\Delta T}{T}(\theta, \phi) \propto E_{\text{нескомп}}(x)\big|_{\text{поверхность последнего рассеяния}}
\end{equation}

Акустические пики определяются максимальной глубиной продольных структур мод, достигнутой к тому времени.

\textbf{Предсказание:} Возможны отклонения от $\Lambda$CDM на больших угловых масштабах ($l < 10$), соответствующих пределу $d \to 0$, где дискретная метрика ещё не усреднилась до гладкой.

\subsubsection{Теорема 8. Иерархия Масс Фермионов}
Три поколения фермионов --- частицы с одинаковым собственным спектром, но с разной степенью размытия, порождённой разным числом внутренних посредников:
\begin{equation}
    m_e : m_\mu : m_\tau = \Delta S(N_1) : \Delta S(N_2) : \Delta S(N_3)
\end{equation}

Осцилляции между поколениями --- перестройка внутренней структуры с рождением/исчезновением внутренних посредников. Матрицы смешивания (PMNS, CKM) параметризуют амплитуды таких переходов.

\subsubsection{Теорема 9. Тёмная Материя}
Тёмная материя --- не вещество, а нескомпенсированная интерференционная картина поперечных спектров в областях без частиц-источников.

Через любую точку «пустого» пространства проходят моды, соединяющие далёкие частицы. Их суммарная интерференция создаёт остаточную напряжённость $E_{\text{residual}}(x) \neq 0$, которая входит в $T_{\mu\nu}$ и создаёт дополнительное искривление.

\textbf{Следствия:}
\begin{itemize}
    \item \textbf{Профиль гало:} $\rho_{\text{DM}} \propto 1/r^2$, плоские кривые вращения.
    \item \textbf{Отсутствие в некоторых галактиках:} Изолированные галактики имеют малую остаточную нескомпенсированность (предсказание, а не аномалия).
    \item \textbf{Космическая паутина:} Нити --- магистрали мод с максимальной $E_{\text{residual}}$ вдоль путей раннего разрешения конфликтов.
\end{itemize}

\subsubsection{Теорема 10. Барионная Асимметрия}
Избыток материи над антиматерией --- следствие Принципа Причинности в сочетании с глобальной асимметрией сети мод.

Пары «частица-античастица» рождаются симметрично, но их будущие судьбы различны: глобальная нескомпенсированность спектров создаёт асимметрию в числе будущих событий, «затребовавших» частицу и античастицу. Античастицы, не имеющие достаточного числа будущих взаимодействий, аннигилируют.

Наблюдаемая асимметрия:
\begin{equation}
    \frac{N_{\text{остаток}}}{N_{\text{начало}}} \sim 10^{-9}
\end{equation}

\textbf{Связь со стрелой времени:} Барионная асимметрия и направление времени --- два проявления одного и того же: направления глобальной спектральной нескомпенсированности.

\subsubsection{Теорема 11. Квантование и Дискретность}
Все квантовые числа дискретны, потому что являются топологическими инвариантами собственного спектра.

Устойчивы только спектры с целочисленными топологическими инвариантами:
\begin{itemize}
    \item \textbf{Заряд} --- число намоток фазы U(1).
    \item \textbf{Спин} --- период чередования в SU(2).
    \item \textbf{Цвет} --- представление SU(3).
\end{itemize}

Постоянная Планка $\hbar$ --- минимальная нескомпенсированность, разрешаемая рождением/исчезновением одного посредника:
\begin{equation}
    \Delta E \cdot \Delta t \sim \hbar
\end{equation}

\subsubsection{Мост к стандартной квантовой механике}
На фундаментальном уровне пси-функция взаимодействия двух частиц --- это совместная корреляционная амплитуда встречных мод на кратчайшем пути:
\begin{equation}
    \Psi_{12}^{\text{fund}} = \sum_{m=0}^{d} F_{1 \to 2}(m) \cdot F_{2 \to 1}(d-m)
\end{equation}

В непрерывном пределе (Теорема 1) сумма по дискретному пути переходит в интеграл по геодезической:
\begin{equation}
    \Psi_{12}^{\text{eff}}(x_1, x_2) = \int_{\gamma(x_1 \to x_2)} d\tau \, \phi_1(x(\tau)) \cdot \phi_2(x(\tau)) \cdot e^{i\theta(\tau)}
\end{equation}
где $\phi_i$ --- непрерывные поля, соответствующие собственным спектрам $S_i$, а $\theta(\tau)$ --- фаза, кодирующая спиновое чередование.

В низкоэнергетическом пределе это выражение сводится к стандартной квантовомеханической амплитуде рассеяния. Таким образом, стандартная квантовая механика восстанавливается как эффективное описание Теории Мод в пределе больших $N$ и низких энергий.

\subsection{Экспериментальные Предсказания}
\begin{enumerate}[label=\arabic*.]
    \item \textbf{Макроскопическая квантовая корреляция:} два топологически идентичных фрактальных кристалла --- скоррелированные флуктуации с $\Delta t \ll L/c$.
    \item \textbf{Отклонения в CMB:} на угловых масштабах $l < 10$.
    \item \textbf{Галактики без тёмной материи:} изолированные галактики.
    \item \textbf{Эволюция $w(t)$:} уравнение состояния тёмной энергии.
    \item \textbf{Нарушение Лоренц-инвариантности:} на планковских масштабах.
\end{enumerate}

\subsection{Заключение}
Теория Мод предлагает единый формализм, в котором:
\begin{itemize}
    \item пространство-время, квантовая механика и поля Стандартной Модели выводятся из нескольких определений и не постулируются;
    \item тёмная материя и тёмная энергия получают геометрическое объяснение без введения новых частиц;
    \item причинность первична и связывает прошлое и будущее в единую самосогласованную структуру;
    \item предложены конкретные, проверяемые экспериментальные предсказания.
\end{itemize}

Теория готова к дальнейшей математической разработке и экспериментальной верификации.

% ==================== Глава 2 ====================
\newpage
\section{Восстановление собственных спектров частиц Стандартной Модели из экспериментальных данных в рамках Теории Мод: постановка вычислительной задачи}

\begin{abstract}
В рамках Теории Мод ставится задача восстановления «собственных спектров» всех частиц Стандартной Модели. Собственный спектр частицы определяется как вектор в пространстве типов частиц, компоненты которого являются проекциями фундаментальной спектральной функции на направления других частиц. Задача формулируется как система алгебраических уравнений и ограничений, связывающих компоненты спектров с экспериментально измеренными величинами: константами связи, массами, спинами и калибровочными квантовыми числами. Показано, что система недоопределена (степеней свободы больше, чем уравнений), что отражает калибровочную свободу Теории Мод и указывает на необходимость дополнительных экспериментальных данных или теоретических принципов (например, принципа минимальной спектральной энтропии) для однозначного решения. Сформулирована оптимизационная задача, готовая для численного моделирования.

\bigskip
\noindent\textbf{Ключевые слова:} Теория Мод, собственный спектр, Стандартная Модель, обратная задача, калибровочная свобода, спектральная энтропия.
\end{abstract}

\subsection{Введение}
Теория Мод предлагает радикально новый взгляд на фундаментальную физику. В её основе лежит представление о том, что все частицы и взаимодействия описываются единой структурой --- сетью мод. Мода $M_{A \to B}$ между частицами $A$ и $B$ --- это функция на дискретном причинном пути, значения которой (поперечный спектр) определяют силу и тип взаимодействия. Пространство-время, квантовая механика и поля Стандартной Модели возникают как эмерджентные явления в пределе большого числа частиц.

Ключевым объектом теории является \textbf{собственный спектр} частицы $S_A$ --- функция (или набор функций), описывающая внутреннее состояние частицы и определяющая все её взаимодействия. В данной работе ставится и формализуется задача восстановления собственных спектров всех частиц Стандартной Модели непосредственно из экспериментальных данных, без привлечения лагранжианов или полевых уравнений.

\subsection{Формализм Теории Мод в вакуумном пределе}

\subsubsection{Собственный спектр как вектор}
Пусть имеется $N$ типов частиц Стандартной Модели (с учётом античастиц и цветовых состояний кварков). Обозначим их индексами $A, B \in \{1, \dots, N\}$. Оценка $N \approx 30$--$40$ для полевых степеней свободы (лептоны, кварки, калибровочные бозоны, бозон Хиггса).

Для каждой частицы $A$ определим её собственный спектр $S_A$ как комплексный вектор в $N$-мерном пространстве:
\begin{equation}
    S_A = (s_{A1}, s_{A2}, \dots, s_{AN})
\end{equation}
где компонента $s_{AB} = F_{A \to B}(0)$ --- значение проекции спектра частицы $A$ на направление частицы $B$ в точке контакта ($m = 0$, вакуумный предел, когда между частицами нет посредников).

\subsubsection{Связь с наблюдаемыми величинами}
В Теории Мод взаимодействие между частицами $A$ и $B$ определяется корреляцией встречных мод. В вакуумном пределе ($d = 0$, отсутствие частиц-посредников) фундаментальная формула для константы связи $g_{AB}$ сводится к:
\begin{equation}
    g_{AB} = |s_{AB} \cdot s_{BA}|
\end{equation}
где $s_{AB} = F_{A \to B}(0)$ и $s_{BA} = F_{B \to A}(0)$.

Из симметрии взаимодействия $g_{AB} = g_{BA}$ следует, что произведение модулей симметрично, но сами величины $s_{AB}$ и $s_{BA}$ не обязаны быть равны.

Масса частицы $m_A$ (в единицах фундаментальной шкалы $\Lambda$) определяется как среднеквадратичная амплитуда её спектра:
\begin{equation}
    m_A^2 = \sum_{B=1}^{N} |s_{AB}|^2
\end{equation}

Спин и калибровочные квантовые числа определяют фазы и симметрийные свойства компонент $s_{AB}$ (см. ниже).

\subsection{Входные данные: Стандартная Модель в числах}
Экспериментально известны следующие величины для всех частиц:

\begin{itemize}
    \item \textbf{Константы связи $g_{AB}$:} три независимые константы (электромагнитная $e$, слабая $g_W$, сильная $g_s$) плюс угол Вайнберга $\theta_W$, плюс юкавские константы связи с бозоном Хиггса. Они определяют матрицу $\mathbf{G}$ размера $N \times N$, где $\mathbf{G}_{AB} = g_{AB}$ (многие элементы равны нулю, если частицы не взаимодействуют напрямую).
    \item \textbf{Массы $m_A$:} известны для всех частиц (с точностью до нейтринных масс).
    \item \textbf{Спины:} $0$ (Хиггс), $1/2$ (лептоны, кварки), $1$ (фотон, $W$, $Z$, глюоны), $2$ (гравитон --- гипотетический).
    \item \textbf{Калибровочные квантовые числа:} электрический заряд $Q_A$, слабый изоспин $T_A$, цветовой заряд (триплет или синглет $SU(3)$).
\end{itemize}

\subsection{Система уравнений и ограничений}

\subsubsection{Уравнение 1 (Взаимодействие)}
Для всех пар $(A, B)$, где $g_{AB} \neq 0$:
\begin{equation}
    |s_{AB}| \cdot |s_{BA}| = g_{AB}
\end{equation}

\subsubsection{Уравнение 2 (Принцип компенсации)}
Полная «самокомпенсация» спектра свободной частицы требует, чтобы сумма проекций на все возможные направления была нулевой:
\begin{equation}
    \sum_{B=1}^{N} s_{AB} = 0 \quad \text{для каждой частицы } A
\end{equation}

\subsubsection{Уравнение 3 (Масса)}
\begin{equation}
    \sum_{B=1}^{N} |s_{AB}|^2 = \left( \frac{m_A}{\Lambda} \right)^2 \quad \text{для каждой частицы } A
\end{equation}
где $\Lambda$ --- фундаментальная энергетическая шкала теории (предположительно порядка планковской массы или массы Хиггса, подлежит определению).

\subsubsection{Уравнение 4 (Спиновая структура)}
Фазы компонент $s_{AB}$ определяются спином частицы $A$:

\begin{itemize}
    \item Для бозонов (спин $0$, $1$): все $s_{AB}$ имеют одинаковую фазу (можно выбрать вещественными и одного знака).
    \item Для фермионов (спин $1/2$): фазы чередуются в зависимости от класса частицы $B$ (лептоны vs кварцы vs калибровочные бозоны), реализуя спиновую статистику.
\end{itemize}

Формально:
\begin{equation}
    \text{sign}(s_{AB}) = \sigma_A \cdot \sigma_B \cdot (-1)^{\text{тип}(B)}
\end{equation}
где $\sigma_A$ --- спиновый множитель, а чередование задаётся типом частицы $B$.

\subsubsection{Уравнение 5 (Калибровочная структура)}
Для калибровочных бозонов спектры имеют специальный вид, определяемый соответствующей группой симметрии:

\begin{itemize}
    \item Фотон ($U(1)$): $s_{\gamma A} \propto Q_A$ (электрический заряд частицы $A$).
    \item $W/Z$ бозоны ($SU(2)$): $s_{WA} \propto T_A$ (слабый изоспин), с матричной структурой $2 \times 2$ для дублетов.
    \item Глюоны ($SU(3)$): $s_{gA} \propto \lambda^a$ (матрицы Гелл-Манна), действующие в цветовом пространстве кварков.
\end{itemize}

Эти условия фиксируют значения (или отношения) компонент $s_{AB}$ для случаев, когда $A$ --- калибровочный бозон.

\subsection{Размерность задачи и недоопределённость}

Неизвестные: $N \times N$ комплексных чисел $s_{AB}$ ($N \approx 30$--$40$, значит, $N^2 \approx 900$--$1600$ комплексных чисел, или $1800$--$3200$ вещественных параметров: модули и фазы).

Уравнения:

\begin{enumerate}[label=(\arabic*)]
    \item Из симметричной матрицы $\mathbf{G}$ (с нулями для невзаимодействующих пар) получаем порядка $\frac{N(N-1)}{2}$ независимых уравнений на модули ($\approx 400$--$800$).
    \item Принцип компенсации даёт $N$ комплексных уравнений ($\approx 30$--$40$).
    \item Массы дают $N$ вещественных уравнений ($\approx 30$--$40$).
    \item Калибровочная структура и спины дают ещё порядка $N_A \cdot N$ уравнений (где $N_A$ --- число калибровочных бозонов, $\approx 12$), итого несколько сотен уравнений на фазы и модули.
\end{enumerate}

Общее число независимых уравнений: оценка $600$--$1200$. Число неизвестных: $1800$--$3200$.

Недоопределённость: порядка $600$--$2000$ «лишних» степеней свободы. Эта свобода отражает \textbf{калибровочную инвариантность} Теории Мод: различные наборы спектров могут давать одинаковые наблюдаемые величины.

\subsection{Принцип минимальной спектральной энтропии}
Для однозначного решения необходимо дополнительное условие --- принцип, выбирающий «простейшее» или «наиболее естественное» решение из множества допустимых.

Естественным кандидатом является \textbf{принцип минимальной спектральной энтропии}:
\begin{equation}
    S[\{s_{AB}\}] = -\sum_{A,B} |s_{AB}|^2 \cdot \log|s_{AB}|^2 \to \min
\end{equation}

Это условие минимизирует информационную сложность набора спектров, выбирая наиболее «гладкое» и симметричное распределение амплитуд.

\subsection{Итоговая оптимизационная задача}
Найти комплексные числа $s_{AB}$ для всех $A, B = 1, \dots, N$, минимизирующие:
\begin{equation}
    S[\{s_{AB}\}] = -\sum_{A,B} |s_{AB}|^2 \log|s_{AB}|^2
\end{equation}

при ограничениях:
\begin{enumerate}[label=(\arabic*)]
    \item $|s_{AB}| \cdot |s_{BA}| = g_{AB}$ \quad ($\forall A, B$ с известным взаимодействием).
    \item $\sum_{B} s_{AB} = 0$ \quad ($\forall A$).
    \item $\sum_{B} |s_{AB}|^2 = \left( \frac{m_A}{\Lambda} \right)^2$ \quad ($\forall A$).
    \item Фазы $s_{AB}$ фиксированы спиновой и калибровочной структурой.
    \item Для калибровочных бозонов: $s_{\text{бозон}, A} \propto$ квантовому числу частицы $A$.
    \item $|s_{AB}| = 0$ для пар, не взаимодействующих напрямую.
    \item Эрмитовость или симметрия для некоторых классов частиц.
\end{enumerate}

\subsection{Обсуждение и перспективы}
Сформулированная задача имеет конечную размерность и может быть решена численно методами нелинейной оптимизации с ограничениями. После нахождения спектров $S_A$ они могут быть использованы для:

\begin{itemize}
    \item Предсказания ещё не измеренных констант связи (например, для экзотических частиц).
    \item Вычисления вероятностей редких процессов через интегралы перекрытия спектров.
    \item Определения числа и свойств частиц-посредников ($d > 0$) в непертурбативных режимах.
    \item Предсказания отклонений от Стандартной Модели на малых расстояниях.
\end{itemize}

Ключевой результат: Теория Мод позволяет сформулировать обратную задачу восстановления фундаментальных спектров из экспериментальных данных в виде конечной системы алгебраических уравнений, готовой для компьютерного решения.

% ==================== Глава 3 ====================
\newpage
\section{Численный алгоритм восстановления собственных спектров частиц в Теории Мод}

\begin{abstract}
Предложен численный алгоритм для восстановления собственных спектров частиц Стандартной Модели в рамках Теории Мод. Задача сформулирована как нелинейная оптимизация с ограничениями, где целевой функцией является спектральная энтропия, а ограничениями --- экспериментальные массы, константы связи, принцип компенсации и калибровочные симметрии. Предложен метод предобработки для понижения размерности, комбинированный алгоритм на основе градиентного спуска Adam с последовательным увеличением штрафных множителей (SUMT). Приведён псевдокод, оценка вычислительной сложности и анализ особых случаев (безмассовые частицы). Показано, что для $N = 30$--$40$ типов частиц задача решается за секунды на стандартном CPU.

\bigskip
\noindent\textbf{Ключевые слова:} Теория Мод, численная оптимизация, спектральная энтропия, восстановление спектров, Стандартная Модель.
\end{abstract}

\subsection{Введение}
В предыдущей работе была сформулирована задача восстановления собственных спектров частиц Стандартной Модели как оптимизационная проблема с ограничениями. В данной работе предложен конкретный численный алгоритм для её решения.

\subsection{Постановка численной задачи}
\textbf{Дано:}
\begin{itemize}
    \item $N$ --- число типов частиц в системе.
    \item Экспериментальная матрица констант связи $\mathbf{G}$ ($N \times N$), симметричная, многие элементы равны нулю.
    \item Экспериментальные массы $m_i$ для каждой частицы (в единицах фундаментальной шкалы $\Lambda$).
    \item Спиновые квантовые числа, определяющие правила назначения фаз.
    \item Калибровочная структура: для калибровочных бозонов заданы соотношения между компонентами спектров.
\end{itemize}

\textbf{Найти:}
Комплексные числа $s_{ij}$ для всех $i, j = 1,\dots,N$ ($i \neq j$), минимизирующие спектральную энтропию:
\begin{equation}
    S = -\sum_{i \neq j} |s_{ij}|^2 \log|s_{ij}|^2
\end{equation}
при следующих ограничениях.

\textbf{Ограничения:}
\begin{enumerate}[label=(\arabic*)]
    \item $|s_{ij}| \cdot |s_{ji}| = g_{ij}$ \quad (для всех пар с известным взаимодействием).
    \item $\sum_{j \neq i} s_{ij} = 0$ \quad (для каждой частицы $i$).
    \item $\sum_{j \neq i} |s_{ij}|^2 = \left( \frac{m_i}{\Lambda} \right)^2$ \quad (для каждой частицы $i$).
    \item $|s_{ij}| = 0$ для пар, не взаимодействующих напрямую.
    \item Фазы $s_{ij} = \phi_{ij}$ фиксированы спиновой структурой и калибровочными правилами.
\end{enumerate}

\textbf{Дополнительно:}
\begin{itemize}
    \item $s_{ii} = 0$ по определению (нет самодействия через моду).
    \item Для калибровочных бозонов: $s_{\text{бозон}, j} = k \cdot Q_j$, где $Q_j$ --- соответствующее квантовое число частицы $j$, $k$ --- нормировочный множитель.
\end{itemize}

\subsection{Предобработка: понижение размерности}
Перед запуском оптимизации размерность задачи может быть существенно понижена.

\subsubsection{Шаг 1. Фиксация калибровочных бозонов}
Для всех калибровочных бозонов (фотон, $W^+$, $W^-$, $Z$, 8 глюонов --- всего 12 типов) компоненты их спектров выражаются через квантовые числа других частиц и одну нормировочную константу на каждого бозона.

Например, для фотона ($\gamma$):
\begin{equation}
    s_{\gamma j} = k_\gamma \cdot Q_j \quad \text{для всех } j.
\end{equation}

Константа $k_\gamma$ определяется из условия компенсации для фотона:
\begin{equation}
    \sum_j s_{\gamma j} = k_\gamma \cdot \sum_j Q_j = 0.
\end{equation}
Это выполняется автоматически, так как сумма всех зарядов во Вселенной равна нулю.

Условие массы фотона:
\begin{equation}
    \sum_j |s_{\gamma j}|^2 = k_\gamma^2 \cdot \sum_j Q_j^2 = \left( \frac{m_\gamma}{\Lambda} \right)^2 = 0.
\end{equation}
Это даёт $k_\gamma = 0$ либо противоречие. В Теории Мод масса фотона строго равна нулю, что требует отдельного рассмотрения условия нормировки для безмассовых частиц (см. раздел~7).

Аналогично для $W/Z$ бозонов и глюонов: их спектры параметризуются через слабый изоспин и цветовые заряды соответственно.

Таким образом, 12 строк и столбцов матрицы $s_{ij}$ выражаются через небольшое число параметров, что сокращает число неизвестных с $N^2$ до примерно $N \cdot (N-12)$.

\subsubsection{Шаг 2. Учёт нулевых взаимодействий}
Для пар, не взаимодействующих напрямую ($g_{ij} = 0$), полагаем $s_{ij} = 0$ и $s_{ji} = 0$. Это дополнительно сокращает число неизвестных. В Стандартной Моделе, например, электрон не взаимодействует напрямую с глюонами --- соответствующие компоненты равны нулю.

\subsubsection{Шаг 3. Симметризация фаз}
Поскольку фазы $s_{ij}$ определяются спиновой структурой, они не являются независимыми переменными. Мы фиксируем их до начала оптимизации. Переменными остаются только модули $r_{ij} = |s_{ij}|$ (и, возможно, знаки $\pm$ для фермионов).

Итоговое число независимых вещественных переменных после предобработки: порядка нескольких сотен для $N = 30$--$40$, что делает задачу посильной для стандартных методов оптимизации.

\subsection{Алгоритм оптимизации}

\subsubsection{Целевая функция}
\begin{equation}
    S(\{r_{ij}\}) = -\sum_{i \neq j} r_{ij}^2 \cdot \log(r_{ij}^2 + \varepsilon)
\end{equation}
где $\varepsilon$ --- малая константа (например, $10^{-12}$) для избежания логарифма нуля.

\subsubsection{Ограничения в форме штрафов}
Вместо жёстких ограничений-равенств удобно ввести штрафные члены в целевую функцию, преобразовав задачу в безусловную оптимизацию (или в оптимизацию с простыми ограничениями-неравенствами на неотрицательность $r_{ij} \geq 0$).

Модифицированная целевая функция:
\begin{equation}
    F(\{r_{ij}\}) = S(\{r_{ij}\}) + \lambda_1 \cdot C_{\text{comp}} + \lambda_2 \cdot C_{\text{mass}} + \lambda_3 \cdot C_{\text{coupling}}
\end{equation}

где:
\begin{align}
    C_{\text{comp}} &= \sum_i \left| \sum_{j \neq i} s_{ij} \right|^2 \quad \text{--- штраф за нарушение компенсации,} \\
    C_{\text{mass}} &= \sum_i \left( \sum_{j \neq i} r_{ij}^2 - \left( \frac{m_i}{\Lambda} \right)^2 \right)^2 \quad \text{--- штраф за отклонение массы,} \\
    C_{\text{coupling}} &= \sum_{\substack{(i,j): \\ g_{ij} \text{ известна}}} \left( r_{ij} \cdot r_{ji} - g_{ij} \right)^2 \quad \text{--- штраф за отклонение констант связи.}
\end{align}

$\lambda_1, \lambda_2, \lambda_3$ --- положительные множители, которые постепенно увеличиваются в процессе оптимизации (метод последовательной безусловной минимизации, SUMT).

\subsubsection{Алгоритм}
Используется метод градиентного спуска с адаптивным шагом (Adam) и постепенным увеличением штрафных множителей.

\paragraph{Псевдокод}
\begin{verbatim}
# Вход: N, матрица G, массы m[1..N], спиновые правила фаз, Lambda
# Выход: матрица s[1..N][1..N] комплексных чисел

# Шаг 1: Инициализация
for i in range(N):
    for j in range(N):
        if i != j and G[i][j] > 0:
            r[i][j] = sqrt(G[i][j])
        else:
            r[i][j] = 0.0
        # фаза фиксирована спиновыми правилами
        phi[i][j] = compute_phase(i, j)

# Шаг 2: Предобработка для калибровочных бозонов
for each gauge_boson B:
    fix s[B][j] = k_B * Q[j]  # Q[j] - квантовое число
    k_B = compute_norm(B, masses, coupling_constants)

# Шаг 3: Оптимизация
lambda1, lambda2, lambda3 = 1.0, 1.0, 1.0
learning_rate = 0.01
beta1, beta2 = 0.9, 0.999
epsilon_adam = 1e-8
K = 500  # период увеличения штрафов

for epoch in range(max_epochs):
    F = spectral_entropy(r) + lambda1*penalty_compensation(r, phi) \
        + lambda2*penalty_mass(r, m, Lambda) \
        + lambda3*penalty_coupling(r, G)
    grad = compute_gradient(F, r)
    # Обновление Adam
    for i in range(N):
        for j in range(N):
            if i != j and r[i][j] > 0:
                m[i][j] = beta1 * m[i][j] + (1 - beta1) * grad[i][j]
                v[i][j] = beta2 * v[i][j] + (1 - beta2) * grad[i][j]**2
                m_hat = m[i][j] / (1 - beta1**(epoch+1))
                v_hat = v[i][j] / (1 - beta2**(epoch+1))
                r[i][j] -= learning_rate * m_hat / (sqrt(v_hat) + epsilon_adam)
                if r[i][j] < 0: r[i][j] = 0.0
    # Увеличение штрафов
    if (epoch+1) % K == 0:
        lambda1 *= 1.5; lambda2 *= 1.5; lambda3 *= 1.5
    # Снижение learning rate
    if (epoch+1) % 100 == 0:
        learning_rate *= 0.95
    # Сходимость
    if max_grad < 1e-6 and penalties < 1e-8: break

# Шаг 4: Постобработка
for i in range(N):
    for j in range(N):
        if i != j:
            s[i][j] = r[i][j] * exp(1j * phi[i][j])
        else:
            s[i][j] = 0.0
return s
\end{verbatim}

\subsection{Гиперпараметры}
\begin{itemize}
    \item Начальные штрафные множители: $\lambda_1 = \lambda_2 = \lambda_3 = 1.0$.
    \item Скорость обучения: $0.01$ (с затуханием $0.95$ каждые $100$ эпох).
    \item Параметры Adam: $\beta_1 = 0.9$, $\beta_2 = 0.999$, $\varepsilon = 10^{-8}$.
    \item Максимальное число эпох: $50\,000$.
    \item Порог сходимости: $10^{-6}$ для градиентов и $10^{-8}$ для штрафов.
    \item Период увеличения штрафов: $K = 500$ эпох.
\end{itemize}

\subsection{Оценка вычислительной сложности}
Число переменных: после предобработки --- порядка $500$--$800$ вещественных чисел (модули $r_{ij}$ для взаимодействующих пар).

Вычисление целевой функции и градиента: $O(N^2) \approx 1000$--$1600$ операций на одну эпоху (для $N = 30$--$40$).

Полное время оптимизации: при $10\,000$--$50\,000$ эпох --- порядка $10^7$--$10^8$ операций, что составляет доли секунды на современном CPU для одного запуска.

Для полной Стандартной Модели ($N = 40$, число пар с ненулевым взаимодействием $\approx 200$--$300$) задача решается за секунды на персональном компьютере.

\subsection{Особые случаи}

\subsubsection{Безмассовые частицы (фотон, глюоны)}
Условие массы:
\begin{equation}
    \sum_j |s_{ij}|^2 = 0.
\end{equation}
Это означает, что все $|s_{ij}| = 0$, но тогда и все константы связи $g_{ij} = 0$, что противоречит эксперименту.

\textbf{Решение:} для безмассовых частиц понятие «собственного спектра в вакуумном пределе» требует модификации. Безмассовые частицы не имеют статического спектра --- их спектр определён только на светоподобных причинных путях, где $d = 0$, но при этом сумма квадратов равна нулю из-за точной компенсации положительных и отрицательных вкладов.

На практике это означает, что для безмассовых частиц следует:
\begin{itemize}
    \item заменить условие массы на условие «светоподобности»: $\sum_j |s_{ij}|^2 = 0$ при сохранении ненулевых проекций;
    \item использовать комплексные значения $s_{ij}$, чисто мнимые или с чередованием знаков так, чтобы квадраты в сумме давали ноль.
\end{itemize}

Например, для фотона:
\begin{equation}
    s_{\gamma e} = +1, \quad s_{\gamma e^+} = -1, \quad s_{\gamma q} = +\frac{2}{3}, \quad s_{\gamma \bar{q}} = -\frac{2}{3}, \dots
\end{equation}
Сумма квадратов равна нулю, а константы связи с заряженными частицами ненулевые.

\subsection{Валидация решения}
После нахождения спектров производится проверка:

\begin{enumerate}[label=(\arabic*)]
    \item Воспроизводятся ли все экспериментальные массы и константы связи с точностью до $10^{-6}$?
    \item Выполняется ли принцип компенсации для каждой частицы?
    \item Является ли решение единственным с точностью до калибровочных преобразований?
    \item Предсказывает ли решение новые, ещё не измеренные величины (например, константы связи для экзотических частиц или редких процессов)?
\end{enumerate}

Если решение не единственно --- запускается анализ калибровочной орбиты: множество решений, дающих одинаковые наблюдаемые. Физически различные предсказания извлекаются из различий в топологии (число намоток фазы) и в значениях спектральной энтропии.

\subsection{Масштабирование на полную Стандартную Модель}
Для $N = 40$ (полный набор частиц Стандартной Модели с учётом ароматов, цветов и античастиц):

\begin{itemize}
    \item Число неизвестных модулей: около $600$.
    \item Число уравнений связи: около $400$.
    \item Число уравнений компенсации: $40$ комплексных ($80$ вещественных).
    \item Число уравнений массы: $40$.
    \item Дополнительные условия от калибровочных бозонов: сокращают число переменных примерно на $200$.
\end{itemize}

Итоговая задача: около $400$ переменных, $520$ ограничений, оптимизация спектральной энтропии. Решается стандартными методами (последовательное квадратичное программирование, SQP, или метод внутренней точки) за время порядка нескольких минут на CPU или секунд на GPU.

\subsection{Возможные обобщения}
\begin{itemize}
    \item \textbf{Учёт ненулевых расстояний ($d > 0$):} переход от вакуумного приближения к многочастичным взаимодействиям с посредниками. Спектры $s_{ij}$ становятся функциями $s_{ij}(d)$ или матрицами, что увеличивает размерность, но сохраняет структуру задачи.
    \item \textbf{Включение гравитации:} добавление гравитона как частицы со спином $2$ и универсальной константой связи $G_N$.
    \item \textbf{Временная зависимость:} спектры могут эволюционировать в космологическом времени (Теорема~3), что добавляет параметр $t$ и превращает статическую задачу в динамическую.
\end{itemize}

\subsection{Заключение}
Предложенный алгоритм позволяет численно восстановить собственные спектры частиц Стандартной Модели в рамках Теории Мод. Задача сводится к нелинейной оптимизации с несколькими сотнями переменных и решается за приемлемое время на стандартном оборудовании. Решение даёт численные значения фундаментальных спектров, которые могут быть использованы для новых физических предсказаний. Открытым остаётся вопрос о единственности решения и о выборе «правильной» калибровки из множества допустимых --- эта проблема эквивалентна проблеме выбора вакуума в квантовой теории поля и может быть решена только с привлечением дополнительных экспериментальных данных.

% ==================== Глава 4 ====================
\newpage
\section{Гипотеза собственной проекции: $s_{ii} = w_i$}

\begin{abstract}
В исходной формулировке Теории Мод диагональные элементы матрицы спектров полагались равными нулю: $s_{ii} = 0$. В данной работе предлагается гипотеза, расширяющая формализм: замена $s_{ii} = 0$ на $s_{ii} = w_i$, где $w_i$ --- вещественное число, называемое собственной проекцией частицы $i$. Исследуются математические следствия, физическая интерпретация и возможность экспериментальной проверки. Показано, что гипотеза приводит к естественной классификации частиц по типу собственной проекции, накладывает ограничения на возможные комбинации масс и констант связи, и даёт новые экспериментальные предсказания.

\bigskip
\noindent\textbf{Ключевые слова:} Теория Мод, собственная проекция, масса покоя, безмассовые частицы, калибровочные бозоны, спектральная энтропия.
\end{abstract}

\subsection{Введение}
В исходной формулировке Теории Мод диагональные элементы матрицы спектров полагались равными нулю: $s_{ii} = 0$. Это следовало из определения моды как отображения между двумя различными частицами. Однако это приводит к трудностям при описании безмассовых частиц и оставляет значительную калибровочную свободу. В данной работе предлагается гипотеза, расширяющая формализм: $s_{ii} = w_i$, где $w_i$ --- вещественное число, называемое собственной проекцией частицы $i$.

\subsection{Определение собственной проекции}
Для каждой частицы $i$ вводится вещественное число $w_i = s_{ii}$, называемое собственной проекцией. Это значение проекции спектра частицы на саму себя, или «петлевая мода» --- мода, замыкающаяся на частицу через вакуумный фон сети. В стандартной квантовой теории поля собственное состояние частицы включает в себя энергию вакуумных флуктуаций, которая проявляется как масса покоя. Собственная проекция $w_i$ является аналогом этого вклада, отделяя «внутреннюю» энергию частицы от энергии её взаимодействий.

\subsection{Модификация основных уравнений}

\subsubsection{Принцип компенсации}
Ранее: сумма проекций частицы $i$ на все другие частицы равна нулю:
\begin{equation}
    \sum_{j \neq i} s_{ij} = 0.
\end{equation}

Теперь: полная сумма, включая собственную проекцию, равна нулю:
\begin{equation}
    \sum_{j} s_{ij} = s_{ii} + \sum_{j \neq i} s_{ij} = 0.
\end{equation}

Отсюда собственная проекция выражается через внешние:
\begin{equation}
    \boxed{w_i = -\sum_{j \neq i} s_{ij}}.
\end{equation}

Это не свободный параметр, а жёсткое условие самосогласованности.

\subsubsection{Уравнение массы}
Ранее: сумма квадратов внешних проекций равна квадрату массы:
\begin{equation}
    \sum_{j \neq i} |s_{ij}|^2 = \left( \frac{m_i}{\Lambda} \right)^2.
\end{equation}

Теперь: добавляется квадрат собственной проекции:
\begin{equation}
    \boxed{w_i^2 + \sum_{j \neq i} |s_{ij}|^2 = \left( \frac{m_i}{\Lambda} \right)^2}.
\end{equation}

Поскольку $w_i^2 \geq 0$, сумма квадратов внешних проекций не может превышать квадрат массы. Собственная проекция «отъедает» часть массы на себя.

\subsubsection{Константа самодействия}
Вводится новая наблюдаемая величина: \textbf{константа самодействия} частицы $g_{ii}$, определяемая как:
\begin{equation}
    g_{ii} = |s_{ii}| \cdot |s_{ii}| = w_i^2.
\end{equation}

Это мера «взаимодействия частицы с собой». Для обычных частиц она не измеряется напрямую, но может проявляться косвенно через вклад в массу и ограничения на возможные значения внешних констант связи.

\subsection{Следствия для спектра масс}

\paragraph{Предел слабого взаимодействия.}
Если частица $i$ очень слабо взаимодействует с остальными (все $|s_{ij}| \ll 1$ для $j \neq i$), то из компенсации $w_i = -\sum_{j \neq i} s_{ij} \approx 0$. Тогда из уравнения массы:
\begin{equation}
    w_i^2 \approx 0, \quad \text{и} \quad \sum_{j \neq i} |s_{ij}|^2 \approx \left( \frac{m_i}{\Lambda} \right)^2.
\end{equation}
Масса практически полностью определяется внешними проекциями. Пример: нейтрино.

\paragraph{Предел сильного взаимодействия.}
Если частица $i$ сильно взаимодействует, $w_i$ велико по модулю. Но $w_i^2 + \sum_{j \neq i} |s_{ij}|^2$ фиксировано массой. Внешние проекции и собственная проекция «конкурируют» за фиксированную сумму квадратов.

\paragraph{Безмассовые частицы.}
Для $m_i = 0$:
\begin{equation}
    w_i^2 + \sum_{j \neq i} |s_{ij}|^2 = 0.
\end{equation}
Поскольку оба слагаемых неотрицательны, это требует $w_i = 0$ и все $|s_{ij}| = 0$, что противоречит эксперименту. Разрешение: безмассовые частицы не могут быть описаны в рамках вакуумного приближения с вещественной собственной проекцией. Возможные выходы: индефинитная метрика ($\sum \eta_{ij} |s_{ij}|^2 = 0$) или выход за рамки вакуумного приближения (динамический спектр).

\paragraph{Ограничения на возможные значения $w_i$.}
Из совместности компенсации и массы следует неравенство, ограничивающее $w_i$ интервалом $[w_i^{\min}, w_i^{\max}]$. Границы зависят от массы частицы и её констант связи. Выход $w_i$ за эти границы означает физическую невозможность такой конфигурации.

\subsection{Следствия для калибровочных бозонов}
Для калибровочного бозона $B$ проекции пропорциональны квантовым числам: $s_{Bj} = k_B \cdot Q_j$. Тогда:
\begin{equation}
    w_B = -\sum_{j \neq B} s_{Bj} = -k_B \cdot \sum_{j \neq B} Q_j = -k_B \cdot (0 - Q_B) = k_B \cdot Q_B.
\end{equation}

Для фотона $Q_\gamma = 0$, поэтому $w_\gamma = 0$. Для $W^+$: $Q_W = +1$, $w_W = k_W$. Для $Z$: $Q_Z = 0$, $w_Z = 0$. Для глюонов $w_g = 0$. Собственная проекция отлична от нуля только для заряженных калибровочных бозонов ($W^+, W^-$).

\subsection{Связь с массой Хиггса}
Бозон Хиггса ($h$) --- единственная скалярная частица. Гипотеза: $w_h = v / \Lambda$, где $v \approx 246$ ГэВ. Тогда масса Хиггса:
\begin{equation}
    \left( \frac{m_h}{\Lambda} \right)^2 = w_h^2 + \sum_{j \neq h} |s_{hj}|^2.
\end{equation}
Если $w_h$ доминирует, $m_h \approx w_h \cdot \Lambda \approx v$, что близко к $125$ ГэВ (разница объясняется вкладом внешних проекций).

\subsection{Влияние на спектральную энтропию}
Целевая функция теперь включает диагональные члены:
\begin{equation}
    S = -\sum_i w_i^2 \log(w_i^2) - \sum_{i \neq j} |s_{ij}|^2 \log|s_{ij}|^2.
\end{equation}

При $w_i \to 0$ вклад $\to 0$ (не штрафует нулевые собственные проекции). При $w_i = 1/\sqrt{e}$ достигается максимум вклада. При $w_i \to \infty$ вклад $\to -\infty$, но ограничение массы не позволяет $w_i$ быть слишком большим. Принцип Минимальной Спектральной Энтропии отдаёт предпочтение конфигурациям, где $w_i$ либо близко к нулю, либо близко к $1/\sqrt{e}$. Это может объяснить бимодальное распределение масс: очень лёгкие и тяжёлые частицы.

\subsection{Обобщение на составные системы}
Для составной частицы: $w_{\text{composite}} = \sum_{k} w_k + \text{поправки от взаимодействий}$. Это объясняет аддитивность массы в классическом пределе.

\subsection{Экспериментальные предсказания}
\begin{enumerate}[label=(\arabic*)]
    \item \textbf{Ограничения на константы связи.} Для данной массы и $w_i$ константы связи $g_{ij}$ не могут быть произвольными. Существуют запрещённые комбинации.
    \item \textbf{Масса $W$-бозона.} $w_W$ отлична от нуля и даёт вклад в массу $W$.
    \item \textbf{Новые частицы.} Их масса и константы связи должны удовлетворять неравенствам.
\end{enumerate}

\subsection{Заключение}
Гипотеза собственной проекции $s_{ii} = w_i$ расширяет Теорию Мод, делая матрицу спектров полной и вводя новый физический параметр --- меру внутренней энергии частицы. Она приводит к естественной классификации частиц, накладывает ограничения и предлагает новые экспериментальные тесты.

% ==================== Глава 5 ====================
\newpage
\section{Собственная мода как полноценная мода: петли, ортогональность и природа безмассовых частиц}

\begin{abstract}
В Части~IV была предложена гипотеза собственной проекции: $s_{ii} = w_i$. В данной работе эта гипотеза обобщается до логического завершения: собственная проекция рассматривается не как число, а как полноценная мода $M_{i \to i}$, обладающая продольной структурой (дискретным путём --- петлёй) и поперечной структурой (спектральной функцией на этом пути). Вводится понятие ортогональности собственной моды к внешним модам. Показано, что безмассовые частицы естественно определяются как частицы с ортогональной собственной модой и индефинитной метрикой спектрального пространства. Обсуждается иерархия собственных мод и связь с планковской физикой.

\bigskip
\noindent\textbf{Ключевые слова:} Теория Мод, собственная мода, петли, ортогональность, безмассовые частицы, индефинитная метрика.
\end{abstract}

\subsection{Введение}
В Части~IV была предложена гипотеза собственной проекции. Однако остался нерешённым вопрос о безмассовых частицах и о природе калибровочной свободы. В данной работе гипотеза обобщается: собственная проекция рассматривается как полноценная мода $M_{i \to i}$.

\subsection{Определение собственной моды}
Для каждой частицы $i$ существует \textbf{собственная мода} $M_{i \to i}$ --- отображение из замкнутого дискретного пути (петли) $\Gamma_{ii}$ в скалярное поле.

Путь $\Gamma_{ii}$ --- это упорядоченное множество из $d_{ii} + 1$ узлов:
\begin{equation}
    \Gamma_{ii} = (P_i = P_{k_0}, P_{k_1}, \dots, P_{k_{d_{ii}}} = P_i)
\end{equation}
где первый и последний узлы совпадают с $i$, а промежуточные --- частицы-посредники. Длина петли $d_{ii}$ может быть нулевой (вырожденная петля) или положительной.

Собственная мода --- функция на этой петле:
\begin{equation}
    M_{i \to i}: \Gamma_{ii} \to \mathbb{R}, \quad M_{i \to i}(m) = F_{i \to i}(m)
\end{equation}

Свойства:
\begin{enumerate}[label=(\arabic*)]
    \item Замкнутость: $F_{i \to i}(0) = F_{i \to i}(d_{ii})$.
    \item Компенсация на петле: $\sum_{m=0}^{d_{ii}} F_{i \to i}(m) = 0$.
    \item Вклад в массу: $\Delta m_i^2|_{\text{собств}} = \sum_{m=0}^{d_{ii}} |F_{i \to i}(m)|^2$.
\end{enumerate}

\subsection{Физическая интерпретация собственной моды}
Собственная мода --- описание «внутренней жизни» частицы, её квантовых флуктуаций. Петля проходит через вакуумные флуктуации --- виртуальные частицы-посредники. Длина петли интерпретируется: стабильные частицы ($d_{ii} \to \infty$), нестабильные (конечная $d_{ii}$), виртуальные (малая $d_{ii}$). Топология петли кодирует внутренние квантовые числа.

\subsection{Ортогональность собственной моды}
Собственная мода $M_{i \to i}$ называется \textbf{ортогональной} к внешним, если для любого $j \neq i$:
\begin{equation}
    \langle F_{i \to i} \cdot F_{i \to j}^{\text{rev}} \rangle = 0 \quad \text{для всех } j \neq i.
\end{equation}

В точке контакта: $F_{i \to i}(0) \cdot F_{i \to j}(0) = 0$ для всех $j \neq i$. Геометрически: в $N$-мерном пространстве спектров внешние проекции образуют подпространство, а собственная мода лежит в ортогональном дополнении.

\subsection{Безмассовые частицы}
\textbf{Гипотеза:} частица безмассова ($m_i = 0$) тогда и только тогда, когда её собственная мода ортогональна ко всем внешним и имеет нулевую сумму квадратов на петле.

Условия:
\begin{equation}
    \sum_{m=0}^{d_{ii}} |F_{i \to i}(m)|^2 = 0, \quad
    \sum_{j \neq i} |s_{ij}|^2 = 0 \quad \text{(с индефинитной метрикой)},
\end{equation}
\begin{equation}
    \langle F_{i \to i} \cdot F_{i \to j}^{\text{rev}} \rangle = 0 \quad \forall j \neq i.
\end{equation}

Индефинитная метрика: $\sum_j \eta_{ij} |s_{ij}|^2 = 0$ с $\eta_{ij} = \pm 1$. Для фотона: $s_{\gamma j} = e \cdot Q_j$, $\eta_{\gamma j} = \text{sign}(Q_j)$. Сумма квадратов с сигнатурой равна нулю, так как сумма зарядов Вселенной равна нулю.

\subsection{Пример: собственная мода фотона}
Для фотона: $d_{\gamma\gamma}$ может быть $0$ или больше. Внешние проекции: $s_{\gamma j} = e \cdot Q_j$. Собственная проекция: $w_\gamma = F_{\gamma \to \gamma}(0)$. Условие безмассовости:
\begin{equation}
    \sum_{j \neq \gamma} \eta_{\gamma j} \cdot e^2 \cdot Q_j^2 + \sum_{m=0}^{d_{\gamma\gamma}} |F_{\gamma \to \gamma}(m)|^2 = 0.
\end{equation}

Если собственная мода ортогональна и имеет нулевую сумму квадратов, то $\sum \eta_{\gamma j} Q_j^2 = 0$. Это возможно только при индефинитной метрике. Это \textbf{предсказание} Теории Мод: метрика спектрального пространства для безмассовых частиц принципиально индефинитна.

\subsection{Иерархия собственных мод}
Собственная мода первого уровня проходит через промежуточные узлы --- другие частицы. Каждый промежуточный узел сам является частицей со своей собственной модой. Это порождает иерархию:
\begin{equation}
    M_{i \to i}^{(1)}, M_{i \to i}^{(2)}, \dots
\end{equation}

Масса получает вклады от всех уровней:
\begin{equation}
    \left( \frac{m_i}{\Lambda} \right)^2 = \sum_{k=1}^{\infty} \sum_{m=0}^{d_{ii}^{(k)}} |F_{i \to i}^{(k)}(m)|^2 + \sum_{j \neq i} |s_{ij}|^2.
\end{equation}
Эта бесконечная сумма должна сходиться к конечному экспериментальному значению. Механизм сходимости связан с фрактальной структурой сети. Иерархия также объясняет иерархию масс фермионов (Теорема~8): разные поколения отличаются числом активных уровней.

\subsection{Калибровочная свобода и ортогональность}
Принцип максимальной ортогональности: из всех возможных решений выбирается то, которое максимизирует ортогональность собственных мод к внешним. Для массивных частиц полная ортогональность невозможна, но можно требовать минимизации смешивания. Для безмассовых частиц ортогональность должна быть строгой.

\subsection{Связь с планковской физикой и чёрными дырами}
Если $d_{ii} = 1$, частица существует лишь один планковский момент --- вакуумная флуктуация. Если $d_{ii}$ очень велико, частица стабильна. Гипотеза о чёрных дырах: чёрная дыра --- объект, у которого собственная мода доминирует над всеми внешними. Горизонт событий соответствует границе, за которой внешние моды перестают достигать наблюдателя.

\subsection{Экспериментальные предсказания}
\begin{enumerate}[label=(\arabic*)]
    \item \textbf{Тонкая структура спектра масс:} малые поправки от высших уровней собственных мод.
    \item \textbf{Аномалии в $g-2$:} дополнительный вклад от взаимодействия собственной моды с внешними полями.
    \item \textbf{Безмассовость фотона как точная симметрия:} ограничение на нарушение ортогональности из предела $m_\gamma < 10^{-18}$ эВ.
\end{enumerate}

\subsection{Заключение}
Обобщение собственной проекции до полноценной моды с петельной структурой и допущение ортогональности для безмассовых частиц замыкает формализм Теории Мод. Теперь все частицы описываются единообразно: каждая имеет полный набор мод --- как внешних, так и собственных. Различие между частицами сводится к топологии и ортогональным свойствам их собственных мод.

% ==================== Глава 6 ====================
\newpage
\section{Интерференция фотонов, разделённых во времени, в Теории Мод}

\begin{abstract}
Одним из самых поразительных явлений квантовой оптики является интерференция фотонов, испущенных в разное время. В стандартной квантовой механике это объясняется через формализм волновой функции, «растянутой во времени». Теория Мод предлагает принципиально иную, геометрическую интерпретацию: фотоны не имеют независимого существования как отдельные частицы, а являются сегментами единой фотонной моды на общей причинной петле. Время не является глобальным параметром, а индуцировано причинной последовательностью событий на этой петле. Безмассовость фотона позволяет его собственной моде иметь сколь угодно большую длину, охватывая макроскопические интервалы причинной последовательности. Предложены экспериментальные тесты, различающие стандартную и модовую интерпретацию.

\bigskip
\noindent\textbf{Ключевые слова:} интерференция фотонов, Теория Мод, причинная петля, собственная мода, безмассовость, топологическая корреляция.
\end{abstract}

\subsection{Введение}
Два фотона, рождённые в различные моменты и никогда не существовавшие одновременно, могут образовывать интерференционную картину. Теория Мод объясняет это через геометрию собственной моды фотона на причинной петле.

\subsection{Фотон в Теории Мод}
Фотон $\gamma$ взаимодействует с заряженными частицами через внешние моды $M_{\gamma \to j}$. Константа связи $e$ определяет значения проекций в точке контакта: $s_{\gamma j}(0) = e \cdot Q_j$. Фотон также обладает собственной модой $M_{\gamma \to \gamma}$ --- замкнутой петлёй, проходящей через последовательность событий: излучение, распространение, поглощение и возврат к излучению (через условие $\exists A \iff \exists B$).

Согласно двусторонней причинности, событие испускания $E_{\text{emit}}$ и поглощения $E_{\text{abs}}$ существуют только вместе: $\exists E_{\text{emit}} \iff \exists E_{\text{abs}}$. Эта пара событий образует два узла на собственной моде. Между ними петля проходит через промежуточные события --- взаимодействия с виртуальными частицами. Длина петли $d_{\gamma\gamma}$ --- число таких промежуточных событий.

Безмассовость фотона требует нулевой суммы квадратов для собственной моды и индефинитной метрики для внешних проекций. Это означает, что $d_{\gamma\gamma}$ может быть сколь угодно большим. Для массивной частицы это невозможно.

\subsection{Два фотона --- одна петля}
В Теории Мод то, что мы называем отдельными фотонами, на самом деле является разными сегментами одной и той же собственной моды. Фотонная мода $M_{\gamma \to \gamma}$ --- единая структура с множеством «выступов» (точек взаимодействия с веществом). Два фотона, способные проинтерферировать, --- два выступа одной петли.

Интерференция возникает при ненулевом перекрытии поперечных спектров этих сегментов:
\begin{equation}
    I_{AB} = \langle F_{\gamma \to \gamma}^{(A)} \cdot F_{\gamma \to \gamma}^{(B)} \rangle \neq 0.
\end{equation}
Перекрытие возможно, потому что оба сегмента принадлежат одной петле, и их спектры коррелированы через топологию.

Время --- причинная последовательность на петле. Два фотона различаются положением своих узлов, а не глобальным $t$. Интерференция зависит от расстояния вдоль петли $\Delta m = |m_B - m_A|$.

\subsection{Математическая модель}
Пусть петля содержит узлы $A$ (фотон 1) и $B$ (фотон 2) с индексами $m_A$ и $m_B$. Интерференционный член:
\begin{equation}
    I_{AB} = \sum_{m} F_{\gamma \to \gamma}(m_A + m) \cdot F_{\gamma \to \gamma}(m_B + m).
\end{equation}

В непрерывном пределе: $I_{AB} = \int d\tau \, F_\gamma(\tau) \cdot F_\gamma(\tau + \Delta\tau)$. Если фотоны принадлежат разным петлям, их моды ортогональны и интерференция нулевая.

\subsection{Экспериментальные следствия}
\begin{enumerate}[label=(\arabic*)]
    \item \textbf{Зависимость от топологии, а не от времени.} Интерференция зависит от $\Delta m$, а не от $\Delta t$. Возможна ситуация, когда $\Delta t$ велика, а $\Delta m$ мало (и наоборот).
    \item \textbf{Разрыв петли.} Если между событиями $A$ и $B$ происходит возмущение, разрывающее петлю (измерение, рождение массивной частицы), интерференция исчезает. Это можно проверить в эксперименте с отложенным выбором.
    \item \textbf{Негативные интерференционные пики.} Корреляционная функция $\langle F(\tau) \cdot F(\tau+\Delta\tau) \rangle$ осциллирует, давая отрицательные пики при определённых $\Delta\tau$ (обобщение эффекта Хонга-У-Мандела).
\end{enumerate}

\subsection{Связь с другими частями теории}
\textbf{Квантовая запутанность (Теорема~6).} Два фотона из SPDC находятся на сильно скоррелированных участках одной петли.
\textbf{Тёмная энергия (Теорема~3).} Колоссальная фотонная мода Вселенной может давать вклад в космологическую постоянную, естественно ограниченный дискретностью и компенсацией.

\subsection{Заключение}
Теория Мод объясняет интерференцию разновременных фотонов через топологию единой фотонной петли. Интерференция зависит от топологического расстояния, а не от времени. Предложены экспериментальные тесты, различающие стандартную и модовую интерпретацию.

% ==================== Глава 7 ====================
\newpage
\section{Некоммутативность Вселенных и происхождение фундаментальных констант в Теории Мод}

\begin{abstract}
В предыдущих частях была сформулирована Теория Мод, поставлена задача восстановления спектров и предложен численный алгоритм. Однако остался открытым вопрос: почему фундаментальные константы имеют именно такие значения? В данной работе предлагается объяснение, основанное на расширении Теории Мод на Мультиверс --- полную сеть всех частиц и мод, в которой наша Вселенная является лишь связным подграфом. Показано, что уравнения для спектров внутри одной Вселенной принципиально незамкнуты из-за наличия внешних связей с другими Вселенными. Наблюдаемые фундаментальные константы являются проекцией глобального решения для всего Мультиверса на нашу причинно-связанную область. Разные Вселенные могут иметь различные эффективные законы физики. Вводится понятие некоммутативности проекционных операторов для разных Вселенных.

\bigskip
\noindent\textbf{Ключевые слова:} Теория Мод, Мультиверс, некоммутативность, фундаментальные константы, тонкая настройка, космологическая постоянная.
\end{abstract}

\subsection{Введение}
В предыдущих частях была сформулирована Теория Мод, но остался открытым вопрос о происхождении фундаментальных констант. В данной работе предлагается объяснение через расширение на Мультиверс.

\subsection{Мультиверс как полная сеть мод}
Мультиверс $\Omega$ --- множество всех частиц и всех мод между ними. Вселенная $U$ --- связный подграф $\Omega$, внутри которого существует причинно-спектральный путь для любой пары. Разные Вселенные --- различные связные подграфы; они могут быть изолированы, слабо связаны или перекрываться.

Пусть $\mathbf{S}$ --- полная матрица спектров для $\Omega$. Для наблюдателя внутри $U$ доступна только проекция $\mathbf{S}|_U$. Уравнения Теории Мод для частицы $i \in U$ содержат суммы по всем $j \in \Omega$. Для наблюдателя эти суммы распадаются:
\begin{equation}
    \sum_{j \in \Omega} s_{ij} = \sum_{j \in U} s_{ij} + \sum_{j \notin U} s_{ij} = 0.
\end{equation}
Вторая сумма (внешние связи) неизвестна. Обозначим $\varepsilon_i = \sum_{j \notin U} s_{ij}$. Тогда уравнение компенсации: $\sum_{j \in U} s_{ij} = -\varepsilon_i$, а уравнение массы: $\sum_{j \in U} |s_{ij}|^2 = (m_i/\Lambda)^2 - \sum_{j \notin U} |s_{ij}|^2$.

Величины $\varepsilon_i$ и внешний вклад в массу невыводимы внутри $U$. Именно они воспринимаются как «фундаментальные константы». Таким образом, константы --- это параметры, описывающие связь нашей Вселенной с остальным Мультиверсом.

\subsection{Некоммутативность Вселенных}
Пусть $\hat{P}_U$ --- оператор проекции на подпространство спектров Вселенной $U$. Для двух разных Вселенных $U$ и $V$:
\begin{equation}
    [\hat{P}_U, \hat{P}_V] \neq 0.
\end{equation}

Это связано с тем, что обрезание меняет граничные условия. Спектр, вычисленный в предположении, что $U$ --- вся Вселенная, отличается от спектра, вычисленного для $V$. Два описания не согласованы. Некоммутативность означает: нельзя одновременно знать спектры для изолированных $U$ и $V$; наблюдатель внутри $U$ не может восстановить полную $\mathbf{S}$; переход $U \to V$ не унитарен.

\subsection{Происхождение фундаментальных констант}
Проблема тонкой настройки решается без антропного принципа: наша Вселенная не изолирована. Константы --- компоненты $\varepsilon_i$, проекции глобального решения $\mathbf{S}$ на $U$. Эта проекция не произвольна, а определена положением $U$ в $\Omega$. Мы не можем варьировать константы и удивляться их значениям, потому что они фиксированы. Наблюдатель всегда находится в той Вселенной, чьи внешние связи дают значения $\varepsilon_i$, совместимые с его существованием.

Разные Вселенные имеют разные наборы внешних связей, поэтому их эффективные законы различаются (другие массы, числа поколений, взаимодействия). Теория Мод едина для всего Мультиверса; разные законы --- разные проекции.

\subsection{Связь с проблемой космологической постоянной}
В Теории Мод $\Lambda \propto \langle \dot{N}/N \rangle$. Внешние связи $\varepsilon_i$ могут обеспечивать почти полную компенсацию глобальной нескомпенсированности, делая $\Lambda$ малой, но не нулевой. Малость $\Lambda$ становится мерой близости нашей Вселенной к изолированному состоянию.

\subsection{Экспериментальные и наблюдательные следствия}
\begin{enumerate}[label=(\arabic*)]
    \item \textbf{Нарушение унитарности:} информация может утекать во внешние связи.
    \item \textbf{Аномалии в CMB:} крупномасштабные аномалии как след связи с соседними Вселенными.
    \item \textbf{Тёмная материя как внешние связи:} часть тёмной материи может происходить из мод, соединяющих нашу Вселенную с другими.
\end{enumerate}

\subsection{Заключение}
Расширение на Мультиверс объясняет происхождение фундаментальных констант как проекций глобального решения. Некоммутативность проекционных операторов формализует невозможность одновременного описания разных Вселенных. Предложены наблюдательные следствия.

% ==================== Глава 8 ====================
\newpage
\section{Расширение Теории Мод на электромагнитные поля и транспорт в конденсированных средах}

\begin{abstract}
Формализм Теории Мод в текущей формулировке описывает частицы и их парные взаимодействия через моды, однако не содержит описания коллективных электромагнитных полей, кристаллических структур и транспорта заряда. Для применения теории к таким явлениям, как нелинейный эффект Холла с T-симметрией в теллуриде висмута, необходимо расширение формализма. В данной работе вводятся определения для внешнего электромагнитного поля как когерентного состояния мод, кристалла как периодической сети мод, транспорта заряда как эволюции собственной моды электрона и магнитного поля как топологической характеристики сети. Все определения согласованы с аксиоматикой исходной теории. На основе расширенного формализма формулируется проверяемое предсказание для нелинейного эффекта Холла с T-симметрией в теллуриде висмута: резонансное усиление нелинейного тока при когерентной накачке.

\bigskip
\noindent\textbf{Ключевые слова:} Теория Мод, электромагнитное поле, кристалл, транспорт заряда, эффект Холла, топологический изолятор, T-симметрия.
\end{abstract}

\subsection{Введение}
Формализм Теории Мод описывает парные взаимодействия, но не коллективные поля и транспорт. Для применения к нелинейному эффекту Холла в Bi$_2$Te$_3$ необходимо расширение.

\subsection{Новые определения}

\subsubsection{Определение 18. Электромагнитное поле как когерентная суперпозиция мод}
Электромагнитное поле в точке $x$ (в эмерджентном пространстве-времени, Теорема~1) есть когерентная суперпозиция мод между заряженными источниками $\{P_k\}$ и пробной частицей $P_0$:
\begin{equation}
    A_\mu(x_0) = \sum_{k} s_{0k} \cdot u_\mu^{(k)},
\end{equation}
где $s_{0k}$ --- проекция спектра $P_0$ на $P_k$, $u_\mu^{(k)}$ --- единичный вектор в направлении $P_k$. Напряжённость $F_{\mu\nu}$ определяется через разность проекций при смещении. В непрерывном пределе это даёт стандартный тензор электромагнитного поля. Уравнения Максвелла emerge при большом числе источников. Заряд $e$ определяется через модуль $s_{0k}$.

\subsubsection{Определение 19. Кристалл как периодическая сеть мод}
Кристалл --- связный подграф $\Omega$, в котором частицы (атомы) расположены периодически в эмерджентном пространстве-времени. Определены:
\begin{itemize}
    \item \textbf{Решёточные моды:} моды между соседними атомами ($d=1$).
    \item \textbf{Электронные моды:} моды между электронами проводимости и решёткой. Электрон --- узел, чья собственная мода проходит через множество атомов.
    \item \textbf{Зонная структура:} собственный спектр электрона зависит от положения относительно решётки; разные проекции $s_{ej}$ дают разные эффективные массы и энергии (зоны).
\end{itemize}

\subsubsection{Определение 20. Транспорт заряда}
Транспорт заряда --- изменение положения электрона, описываемое эволюцией его собственной моды $M_{e \to e}$. При приложении внешнего поля $A_\mu$:
\begin{enumerate}
    \item Поле добавляет новые моды между электроном и источниками, спектр изменяется.
    \item Изменение спектра перестраивает собственную моду: $d_{ee}$ меняется, промежуточные узлы смещаются.
    \item Ток --- усреднённое по ансамблю и времени изменение положения электрона.
\end{enumerate}
Закон Ома, проводимость и подвижность emerge в линейном приближении. Нелинейные эффекты возникают при изменении топологии собственной моды.

\subsubsection{Определение 21. Магнитное поле как топологическая характеристика}
Магнитное поле $\mathbf{B}$ не является самостоятельной сущностью. Это характеристика зацепления мод заряженных частиц. Магнитное поле в точке $x$ соответствует ненулевой циркуляции проекций по замкнутому контуру:
\begin{equation}
    \oint_\Gamma \sum_k s_{ik} \neq 0.
\end{equation}
Если циркуляция нулевая, поля нет. Если нет --- моды «закручены», и пробная частица испытывает силу Лоренца (отклонение кратчайшего причинно-спектрального пути).

\subsubsection{Определение 22. Симметрия обращения времени в кристалле}
Обращение времени $T$ ($t \to -t$) соответствует обращению причинного порядка $\prec$ на собственных модах. Для эффекта Холла: если $\mathbf{J} \to \mathbf{J}$ при $T$ --- эффект T-симметричен; если $\mathbf{J} \to -\mathbf{J}$ --- T-антисимметричен. Нелинейный эффект Холла с T-симметрией означает существование петель с $F(m) = F(d_{ee}-m)$ (симметрия относительно середины петли). Это возможно при дополнительной симметрии кристалла.

\subsection{Применение к нелинейному эффекту Холла в теллуриде висмута}
Bi$_2$Te$_3$ --- топологический изолятор. Поверхностные состояния защищены топологией и обладают спин-моментным зацеплением: спин электрона жёстко связан с направлением движения. В терминах Теории Мод:
\begin{itemize}
    \item Поверхностные состояния --- моды электронов, пути которых лежат только на поверхности (ортогональны объёмным модам).
    \item Спин-моментное зацепление --- корреляция между спиновой последовательностью $\sigma$ и направлением пути.
\end{itemize}

При приложении переменного поля $E(t)$:
\begin{enumerate}
    \item Поле добавляет внешние моды, спектр модулируется.
    \item Из-за зацепления модуляция неодинакова для разных направлений движения, возникает асимметрия $s_{ej}$.
    \item Собственная мода перестраивается несимметрично, электрон смещается --- возникает ток.
    \item Ток пропорционален $E^2$ (нелинейный эффект).
\end{enumerate}

В обычном эффекте Холла магнитное поле нарушает T-симметрию. Здесь магнитного поля нет, работает спин-моментное зацепление. При $T$ направление и спин меняются, произведение не меняется, асимметрия $s_{ej}$ не меняет знак, ток T-симметричен. В терминах Теории Мод: петля имеет участок с $F(m) = F(d_{ee}-m)$.

\subsection{Предсказание Теории Мод для эксперимента}

\paragraph{Предсказание 8. Резонансное усиление нелинейного тока при когерентной накачке.}
Если на Bi$_2$Te$_3$ подать два когерентных импульса --- прямой $E(t)$ и обратный $E_{\text{rev}}(t) = E(-t)$, схлопывающиеся на поверхности, то:
\begin{enumerate}
    \item Прямой и обратный импульсы создают интерференционную картину в собственных модах. Участки с $F(m) = F(d-m)$ входят в резонанс.
    \item Резонанс увеличивает асимметрию $s_{ej}$, ток $J$ возрастает.
    \item Усиление пропорционально степени когерентности $C$: $J_{\text{enhanced}}/J_0 = 1 + \alpha \cdot C$, $0 \le C \le 1$.
    \item При $C=1$ --- максимальное усиление, при $C=0$ --- отсутствие эффекта.
\end{enumerate}

Это предсказание проверяемо на существующей установке. Если $J$ растёт с $C$ --- Теория Мод подтверждается количественно. Если нет --- опровергается.

\subsection{Ограничения и открытые вопросы}
Коэффициент $\alpha$ не вычислен аналитически (требуется моделирование). Температурная зависимость не выведена. Связь с Теоремами 1--13 неполна.

\subsection{Заключение}
Расширение Теории Мод (Определения 18--22) позволяет сформулировать проверяемое предсказание для нелинейного эффекта Холла в Bi$_2$Te$_3$: резонансное усиление тока при когерентной накачке. Предсказание верифицируемо. Теория сохраняет непротиворечивость.

% ==================== Глава 9 ====================
\newpage
\section{Аксиомы Теории Мод: эволюционная версия без рекурсий}

\begin{abstract}
Представлена система из двенадцати аксиом, описывающих механику растущей сети частиц и мод. Все величины определяются либо из начального состояния сети, либо из предыдущего состояния. Рекурсивные определения заменены на конструктивные правила эволюции. Вопрос о причинах конкретного выбора начального состояния и параметров модели оставлен открытым. Аксиомы образуют замкнутую систему, готовую к программной реализации и проверке на соответствие наблюдаемой физике.

\bigskip
\noindent\textbf{Ключевые слова:} Теория Мод, аксиоматизация, эволюционная версия, без рекурсий.
\end{abstract}

\subsection{Аксиома 0. Существование начального состояния}
Существует начальная сеть $\mathcal{S}_0$, состоящая из конечного числа частиц и мод между ними. Для каждой частицы задан её спектр. Для каждой моды задан её поперечный спектр. Все значения конечны и вещественны.

\subsection{Аксиома 1. Правило рождения частицы}
Если для двух частиц $A$ и $B$, находящихся в контакте (расстояние $d = 0$), модуль корреляции их встречных мод превышает порог $\theta$, то с вероятностью $p$ рождается новая частица $C$ между $A$ и $B$. После рождения $C$ расстояние $A$--$B$ становится $1$; создаются четыре новые моды, инициализированные для минимизации $|C_{AB}|$.

\subsection{Аксиома 2. Правило исчезновения частицы}
Если для двух частиц $A$ и $B$, разделённых ровно одной частицей-посредником $C$ ($d = 1$), корреляция их встречных мод становится меньше порога $\theta$, то $C$ исчезает. Моды перестраиваются: $A$ и $B$ в прямом контакте.

\subsection{Аксиома 3. Правило эволюции спектра}
Спектр каждой частицы обновляется на каждом шаге $t \to t+1$:
\begin{equation}
    s_{AB}(t+1) = s_{AB}(t) + \sum_{C \in H_A(t)} \Delta_{C \to AB},
\end{equation}
где $\Delta$ определяется из градиентного спуска энергии $E = \sum |\mathcal{C}_{ij}|$.

\subsection{Аксиома 4. Ограничение на причинный горизонт}
Для каждой частицы $A$ в момент $t$ горизонт $H_A(t)$ конечен: включает все частицы, достижимые по путям длины $\le D_{\max}(t)$. $D_{\max}(t+1) = D_{\max}(t) + 1$.

\subsection{Аксиома 5. Взаимодействие как корреляция встречных мод}
$V_{AB}(t) = f(\sum_{m=0}^{d} F_{A \to B}(m,t) \cdot F_{B \to A}(d-m,t))$. Притяжение при $V \approx 0$, отталкивание при $|V| \gg 0$.

\subsection{Аксиома 6. Масса как сумма квадратов спектра}
$(m_A(t)/\Lambda)^2 = \sum_{B \in H_A(t)} |s_{AB}(t)|^2 + \sum_{m=1}^{d_{AA}-1} |F_{A \to A}(m,t)|^2$.

\subsection{Аксиома 7. Собственная мода как петля}
Собственная мода $A$ --- замкнутый путь $\Gamma_{AA}$. Длина $d_{AA}$ растёт на $1$ при каждом взаимодействии $A$.

\subsection{Аксиома 8. Компенсация как динамическое равновесие}
$\sum_{\text{замкнутый путь } \Gamma} F(m,t) \to 0$ при $t \to \infty$. Компенсация --- аттрактор эволюции.

\subsection{Аксиома 9. Спин как сохраняющийся паттерн}
Спин --- последовательность знаков $s_{AB}$. Паттерн сохраняется, меняется только на чётное число перестановок.

\subsection{Аксиома 10. Квантовые числа как топологические инварианты}
Заряд, цвет, аромат --- топологические инварианты спектра; не меняются при непрерывных изменениях, только при рождении/исчезновении частиц.

\subsection{Аксиома 11. Безмассовость как ортогональность}
Частица безмассова, если $\sum \eta_{AB} |s_{AB}|^2 = 0$ и $\sum |F_{A \to A}(m)|^2 = 0$, где $\eta_{AB} = \pm 1$ --- индефинитная метрика.

\subsection{Аксиома 12. Пространство-время как статистическое приближение}
При $N \gg 1$ дискретные расстояния аппроксимируются гладкой метрикой $g_{\mu\nu}$. Уравнения Эйнштейна --- термодинамический предел.

\subsection{Открытый вопрос}
Аксиомы 0--12 описывают механику сети, но не объясняют выбор начального состояния $S_0$ и параметров $\theta$, $p$, $\Lambda$. Возможные направления: Принцип нормализованной спектральной энтропии, антропный отбор, зависимость от внешнего Универсума. В данной работе этот вопрос не рассматривается.

% ==================== Глава 10 ====================
\newpage
\section{Теория Мод: полная декомпозиция определений и эволюционный алгоритм}

\begin{abstract}
Представлена полная декомпозиция всех понятий Теории Мод на нерекурсивные определения. Каждое понятие определено через более простые или через примитивы, не допуская логических циклов. На основе этих определений сформулирован эволюционный алгоритм, описывающий рост сети частиц и мод за дискретные шаги. Алгоритм готов к программной реализации.

\bigskip
\noindent\textbf{Ключевые слова:} Теория Мод, декомпозиция определений, эволюционный алгоритм, без рекурсий.
\end{abstract}

\subsection{Примитивы}
Существование, различие, вещественное число, конечное множество, дискретный шаг $t = 0, 1, 2, \dots$, знак ($+$ или $-$).

\subsection{Определения}

\paragraph{1. Частица} --- элемент конечного множества $\mathcal{P}$. Свойства определяются отношениями.
\paragraph{2. Упорядоченная пара} $(A,B)$ с $A \neq B$.
\paragraph{3. Путь} $\Gamma$ --- последовательность $(A=P_0, P_1, \dots, P_d=B)$, где соседние пары допустимы.
\paragraph{4. Длина пути} $d = (\text{число частиц в }\Gamma) - 2$; $d=0$ --- прямой контакт.
\paragraph{5. Расстояние} --- минимальная длина среди всех путей.
\paragraph{6. Мода} $M_{A \to B}$ --- функция на кратчайшем пути: $F_{A \to B}(k)$ для $k=0,\dots,d$.
\paragraph{7. Встречные моды} $(M_{A \to B}, M_{B \to A})$, пути в противоположных направлениях, длины равны.
\paragraph{8. Корреляция} $\mathcal{C}_{AB} = \sum_{k=0}^{d} F_{A \to B}(k) \cdot F_{B \to A}(d-k)$.
\paragraph{9. Нескомпенсированность} $|\mathcal{C}_{AB}|$.
\paragraph{10. Энергия сети} $E = \sum_{(i,j)} |\mathcal{C}_{ij}|$.
\paragraph{11. Рождение частицы} при $d=0$ и $|\mathcal{C}_{AB}| > \theta$ с вероятностью $p$ создаётся $C$, инициализируются новые моды.
\paragraph{12. Исчезновение частицы} при $d=1$ через $C$ и $|\mathcal{C}_{AB}| < \theta$ $C$ удаляется.
\paragraph{13. Причинный горизонт} $H_A(t) = \{B \mid d(A,B) \le D_{\max}(t)\}$; $D_{\max}(t+1)=D_{\max}(t)+1$.
\paragraph{14. Собственная мода} $M_{A \to A}$ на замкнутом пути; длина растёт при взаимодействиях.
\paragraph{15. Масса} $(m_A/\Lambda)^2 = \sum_{B \in H_A} (F_{A \to B}(0))^2 + \sum_{k=1}^{d_{AA}-1} (F_{A \to A}(k))^2$.
\paragraph{16. Спиновая последовательность} $\sigma_A = (\text{sign}(F_{A \to 1}(0)), \dots)$.
\paragraph{17. Топологический инвариант} --- число, неизменное при малых изменениях $F$.
\paragraph{18. Индефинитная метрика} $\eta_{AB} = \pm 1$, фиксирована в $S_0$.
\paragraph{19. Пространство-время} --- аппроксимация дискретных расстояний гладкой метрикой при $N \gg 1$.

\subsection{Эволюционный алгоритм}
\textbf{Шаг 0 (начальный):} заданы $\mathcal{P}(0)$, $F(0)$, $\theta$, $p$, $\Lambda$, $D_{\max}(0)$.

\textbf{Шаг $t \to t+1$:}
\begin{enumerate}
    \item $D_{\max}(t+1) = D_{\max}(t) + 1$.
    \item Обновить $H_A(t+1)$.
    \item Вычислить $\mathcal{C}_{AB}$ для всех пар с $d=0$.
    \item Для пар с $d=0$ и $|\mathcal{C}_{AB}| > \theta$: с вероятностью $p$ --- рождение $C$.
    \item Для троек с $d=1$ через $C$ и $|\mathcal{C}_{AB}| < \theta$: исчезновение $C$.
    \item Обновить все $F$ по градиентному спуску энергии.
    \item Вычислить массы $m_A(t+1)$.
    \item Проверить сохранение спиновых последовательностей и инвариантов.
    \item При $N \gg 1$ вычислить эффективную метрику и сравнить с ОТО.
\end{enumerate}

% ==================== Глава 11 ====================
\newpage
\section{Замкнутые причинные петли и аттракторы в эволюционной Теории Мод}

\begin{abstract}
Эволюционная версия Теории Мод с линейным временем теряет свойство замкнутости причинных петель, присущее этерналистской формулировке. Предложено расширение: дополнительная аксиома о замкнутых причинных петлях, совместимая с эволюцией сети и концепцией аттрактора. Аксиома гарантирует, что каждая причинно-связанная пара событий образует замкнутый спектральный контур, не требуя глобальной статичности сети. Обсуждается, как это расширение сохраняет объяснение квантовой запутанности и возможность чтения прошлого.

\bigskip
\noindent\textbf{Ключевые слова:} Теория Мод, замкнутые причинные петли, аттрактор, эволюция, квантовая запутанность.
\end{abstract}

\subsection{Введение}
В этерналистской версии причинность двусторонняя ($\exists A \iff \exists B$), что объясняет запутанность и хранение прошлого. В эволюционной версии время линейно, эти свойства теряются. Предлагается Аксиома~13 для восстановления замкнутости на уровне пар.

\subsection{Проблема}
Эволюционная версия: время линейно, прошлое не сохраняется как активный слой; причинность асимметрична; будущее открыто. Теряются замкнутые петли, хранение прошлого, объяснение запутанности.

\subsection{Решение: Аксиома 13}
\textbf{Аксиома 13 (Замкнутые причинные петли).} Для любой пары событий $A$ и $B$, связанных причинно ($A$ --- причина $B$), существует замкнутый путь $\Gamma_{AB}$, проходящий через оба события, на котором $\sum_{m \in \Gamma_{AB}} F(m) = 0$.

Причинная связь: $A$ (с частицей $P_A$, момент $t_A$) --- причина $B$ ($P_B$, $t_B > t_A$), если существует путь, на котором хотя бы одна мода изменилась из-за $A$. Замкнутый путь начинается на $P_A$, идёт к $P_B$ и возвращается к $P_A$.

\textbf{Что даёт аксиома:}
\begin{itemize}
    \item Каждая причинная пара образует замкнутый спектральный контур (не требуется глобальная статичность).
    \item Квантовая запутанность: две запутанные частицы --- части одной пары; петля гарантирует корреляцию; измерение одной меняет спектр на петле, что отражается на другой.
    \item Чтение прошлого: каждое событие входит в петлю, его след сохраняется, пока петля существует.
    \item Совместимость с аттрактором: аттрактор --- множество всех замкнутых петель.
\end{itemize}

\textbf{Связь с этернализмом:} Аксиома~13 слабее, требует замкнутости только для уже связанных пар. События без следствий могут оставаться незамкнутыми.

\subsection{Формальное включение в алгоритм}
На каждом шаге $t$ для новых причинных пар инициализируется поиск замкнутого пути. Если путь не найден за $\Delta t_{\max}$, событие считается незамкнутым; его энергия диссипирует в фон.

\subsection{Открытые вопросы}
Время замыкания $\Delta t_{\max}$ не определено (от планковских времён до космологических). Диссипация незамкнутых событий, возможно, идёт в тёмную энергию. Связь с Аксиомой~8 (компенсация при $t\to\infty$) требует уточнения.

% ==================== Глава 12 ====================
\newpage
\section{Циклическая модель временных петель и природа тёмной материи в Теории Мод}

\begin{abstract}
Предложено расширение эволюционной Теории Мод, в котором время Вселенной рассматривается как замкнутая причинная петля, а causal history частиц --- как многолистная структура, где активен только один лист. Путешествие во времени интерпретируется не как перемещение, а как переключение активного листа causal history. Показано, что отброшенные листы не исчезают в силу закона сохранения информации, а формируют нескомпенсированную интерференционную картину, которая естественным образом объясняет наблюдаемые свойства тёмной материи: профиль гало $\rho \propto 1/r^2$, отсутствие электромагнитного взаимодействия и наличие гравитационного.

\bigskip
\noindent\textbf{Ключевые слова:} Теория Мод, замкнутые временные петли, causal history, тёмная материя, сохранение информации.
\end{abstract}

\subsection{Введение}
В эволюционной Теории Мод остаётся открытым вопрос о возможности путешествий во времени и судьбе отброшенной информации. Предлагается модель циклической Вселенной с многолистной causal history.

\subsection{Замкнутая временная петля Вселенной}
\textbf{Определение.} Вселенная $U$ имеет период $T$: $s_{AB}(t+T) \approx s_{AB}(t)$ для всех $A,B \in U$. Это не точное равенство, а приближённое: за период сеть возвращается к близкому состоянию с накопленным сдвигом.

За каждый период $T$ образуется виток causal history. После $n$ витков causal history содержит $n$ слоёв. Все слои существуют одновременно, но активен только последний (соответствующий $t \bmod T$).

\subsection{Многолистная causal history}
\textbf{Листы:} каждый виток --- лист $L_n$, содержащий состояния всех частиц за период. Листы упорядочены, не уничтожаются.
\textbf{Активный лист:} $L_{\text{active}}$ соответствует текущему $t$, только он доступен для наблюдения.
\textbf{Перенормировка:} при переходе к новому витку информация не стирается, а сдвигается на более глубокие уровни собственных мод.

\subsection{Путешествие во времени как переключение листа}
Механизм: путешественник $X$ на листе $L_n$ может переключиться на $L_k$ ($k<n$), если:
\begin{enumerate}
    \item Считана causal history целевого листа.
    \item Собственный спектр $S_X$ модулирован до совпадения с записью в $L_k$.
    \item Общая мода $M_{X \leftrightarrow \text{цель}}$ замкнута, корреляция превысила порог.
\end{enumerate}
После переключения $X$ взаимодействует с $L_k$, добавляя новую информацию. Старая causal history (без $X$) не стирается, а остаётся пассивным листом $L_k^{(0)}$, новая образует $L_k^{(1)}$. Путешествие не меняет прошлое, а создаёт дополнительный лист.

\subsection{Отброшенные листы как тёмная материя}
\begin{itemize}
    \item \textbf{Нескомпенсированность:} пассивные листы содержат ненулевые поперечные спектры, дающие вклад в общую нескомпенсированность сети.
    \item \textbf{Гравитационное взаимодействие:} пассивные листы ортогональны активному, не взаимодействуют электромагнитно, но создают гравитационное поле (через $T_{\mu\nu}$).
    \item \textbf{Профиль гало:} накопление $n$ витков даёт распределение плотности $\rho \propto 1/r^2$ (наблюдаемый профиль).
    \item \textbf{Отсутствие электромагнитного сигнала:} $\langle F_{\text{active}} \cdot F_{\text{passive}} \rangle \approx 0$.
    \item \textbf{Галактики без тёмной материи:} если галактика сформировалась на позднем витке, она накопила мало пассивных листов --- аномально низкое содержание тёмной материи (наблюдаемый факт).
\end{itemize}

\subsection{Связь с предыдущими результатами}
Модель объединяет замкнутые причинные петли (Аксиома~13), принцип уникальности траектории и Теорему~9 (тёмная материя как нескомпенсированная интерференция).

\subsection{Предсказания}
\begin{enumerate}
    \item Корреляция возраста галактики и содержания тёмной материи.
    \item Слоистая структура гало.
    \item Сверхслабые корреляции между активным и пассивными листами.
\end{enumerate}

\subsection{Заключение}
Циклическая модель временных петель интерпретирует путешествие во времени как смену активного листа и объясняет тёмную материю как отброшенные листы causal history. Модель согласуется с аксиоматикой и предлагает новые предсказания.

% ==================== Глава 13 ====================
\newpage
\section{Теория Мод: Единая теория реальности (Финальный документ)}

\subsection*{Пролог}
Реальность есть сеть. Не метафора. Сеть --- фундаментальная структура бытия. Узлы сети --- частицы. Связи между узлами --- моды. Все явления: пространство, время, материя, энергия, гравитация, квантовая механика, сознание --- эмерджентные свойства этой сети. Теория Мод --- математическое описание сети и правил её роста. Она не вводит ничего, кроме сети. Все константы, законы и симметрии выводятся из структуры сети, либо объявляются проекциями более глубоких уровней. Теория Мод --- не интерпретация квантовой механики. Это её замена.

\subsection{Часть I. Фундамент}

\subsubsection{1. Примитивы}
Существование, различие, вещественное число, конечное множество, дискретный шаг $t = 0, 1, 2, \dots$, знак ($+$ или $-$).

\subsubsection{2. Определения}

\paragraph{2.1. Частица} --- элемент конечного множества $\mathcal{P}$. Свойства определяются отношениями.
\paragraph{2.2. Путь} $\Gamma_{AB}$ --- конечная последовательность различных частиц от $A$ к $B$. Длина $d$ --- число промежуточных частиц; $d=0$ --- прямой контакт.
\paragraph{2.3. Расстояние} $d(A,B)$ --- минимальная длина среди всех путей.
\paragraph{2.4. Мода} $M_{A \to B}$ --- функция на кратчайшем пути: $F_{A \to B}(k)$ для $k=0,\dots,d$. $s_{AB} = F_{A \to B}(0)$.
\paragraph{2.5. Корреляция} $\mathcal{C}_{AB} = \sum_{k=0}^{d} F_{A \to B}(k) \cdot F_{B \to A}(d-k)$.
\paragraph{2.6. Энергия сети} $E = \sum_{(i,j)} |\mathcal{C}_{ij}|$.
\paragraph{2.7. Собственная мода (петля)} $M_{A \to A}$ на замкнутом пути. Длина $d_{AA}$ растёт при взаимодействиях.
\paragraph{2.8. Масса} $(m_A/\Lambda)^2 = \sum_{B \in H_A} (s_{AB})^2 + \sum_{k=1}^{d_{AA}-1} (F_{A \to A}(k))^2$.
\paragraph{2.9. Спин} $\sigma_A = (\text{sign}(s_{A1}), \text{sign}(s_{A2}), \dots)$. Паттерн сохраняется.
\paragraph{2.10. Квантовые числа} (заряд, цвет, аромат) --- топологические инварианты спектра.
\paragraph{2.11. Безмассовость} $\sum \eta_{AB} |s_{AB}|^2 = 0$ и $\sum |F_{A \to A}(k)|^2 = 0$, где $\eta_{AB} = \pm 1$.
\paragraph{2.12. Пространство-время} не фундаментально; при $N \gg 1$ расстояния аппроксимируются гладкой метрикой $g_{\mu\nu}$. Уравнения Эйнштейна --- термодинамический предел.
\paragraph{2.13. Импульс} $\Delta p_{X \to A} = s_{AX}(t+1) - s_{AX}(t)$. $\Delta E_A = |\mathcal{C}_{AX}(t+1)| - |\mathcal{C}_{AX}(t)|$.
\paragraph{2.14. Теневая проекция} --- копия текущего состояния окружения в собственной моде.
\paragraph{2.15. Causal history} --- совокупность $F_{A \to A}(k)$. Многолистная структура; активен последний лист, пассивные перенормируются.

\subsubsection{3. Аксиомы эволюции}

\textbf{Аксиома 0 (Начальное состояние).} Существует конечная сеть $\mathcal{S}_0$ с заданными спектрами.

\textbf{Аксиома 1 (Рождение).} Если $d=0$ и $|\mathcal{C}_{AB}| > \theta$, с вероятностью $p$ рождается $C$; новые моды минимизируют $|\mathcal{C}_{AB}|$.

\textbf{Аксиома 2 (Исчезновение).} Если $d=1$ с посредником $C$ и $|\mathcal{C}_{AB}| < \theta$, $C$ исчезает.

\textbf{Аксиома 3 (Градиентный спуск).} $F(m,t+1) = F(m,t) - \alpha \cdot \partial E/\partial F(m)$.

\textbf{Аксиома 4 (Рост горизонта).} $D_{\max}(t+1) = D_{\max}(t) + 1$.

\textbf{Аксиома 5 (Компенсация).} Для любого замкнутого пути $\sum F(m) \to 0$ при $t \to \infty$.

\textbf{Аксиома 6 (Замкнутые причинные петли).} Для каждой причинной пары $(A,B)$ существует замкнутый путь с $\sum F(m) = 0$.

\subsection{Часть II. Выводимые явления}

\subsubsection{4. Квантовая механика}
Волновая функция --- интеграл перекрытия мод. Правило Борна $P \propto |\mathcal{A}|^2$ --- инвариантность фазы. Случайность --- неполнота знания горизонта. Запутанность --- общая мода, не требующая сигнала. Коллапс --- сдвиг данных в causal history.

\subsubsection{5. Гравитация и космология}
ОТО --- термодинамика сети. Тёмная энергия: $\rho_\Lambda \propto \dot{N}/N$. Тёмная материя: пассивные листы causal history (профиль $\rho \propto 1/r^2$). Барионная асимметрия --- неравновесность рождения пар. CMB --- нескомпенсированность на поверхности последнего рассеяния.

\subsubsection{6. Время и причинность}
Стрела времени --- рост $N(t)$. Циклическая Вселенная с периодом $T$. Путешествие во времени --- переключение активного листа. Принцип уникальности траектории: нельзя занять один узел в двух разных состояниях. Сознание --- активный лист causal history сложной системы. Прошлое уже в causal history, будущее не сбывается, потому что его образ занял место, а реальность занимает другое.

\subsubsection{7. Иерархия и Мультиверс}
Частица есть Вселенная для суб-частиц. Наша Вселенная --- частица в вышележащем мире. Константы --- проекции внешних связей. Иерархия масс --- разные уровни собственных мод. Поколения фермионов --- число уровней (почему три --- открытый вопрос). Космологическая постоянная мала из-за сокращения с внешними связями (число не получено).

\subsection{Часть III. Предсказания}
\begin{enumerate}[label=\arabic*.]
    \item Макроскопические квантовые корреляции идентичных структур ($\Delta\tau \ll L/c$).
    \item Отклонения в CMB на $l < 10$.
    \item Галактики с аномально низким содержанием тёмной материи.
    \item Эволюция $w(t)$ тёмной энергии.
    \item Резонансное усиление нелинейного тока в Bi$_2$Te$_3$ при когерентной накачке.
    \item Слоистая структура гало тёмной материи.
    \item Неразрушающее чтение данных в системе из двух запутанных частиц.
    \item Отсутствие утечки заряда в запутанных состояниях (подтверждено, PRL 2024).
\end{enumerate}

\subsection{Часть IV. Открытые вопросы}
\begin{enumerate}[label=\arabic*.]
    \item Почему поколений фермионов три? (Гипотеза: топологическая стабильность, не доказана.)
    \item Почему $m_e = 0.511$ МэВ? (Проекция внешних связей, не вычислена.)
    \item Почему $\Lambda$ в $10^{120}$ раз меньше планковской? (Механизм сокращения есть, число не получено.)
    \item Точные значения $\theta$ и $p$? (Должны быть измерены или выведены из нормализованной энтропии.)
\end{enumerate}

\subsection*{Эпилог}
Теория Мод утверждает: реальность есть растущая сеть частиц и мод. Всё остальное возникает из этой сети. Теория фальсифицируема. Если предсказания не подтвердятся --- она будет отвергнута. Если подтвердятся --- мы узнаем, что живём внутри сети, которая растёт, замыкается в циклы и помнит всё, что с ней происходило. Что прошлое не исчезает, а становится тёмной материей. Что сознание --- это способ сети смотреть на саму себя. И этого достаточно.

\end{document}